\section{Absolute cohomology: the finite-field identification and its arithmetic limits}\label{sec:part-b}
\providecommand{\BW}{\mathrm W}
\providecommand{\BM}{\mathrm M}
\providecommand{\Bcts}{\mathrm{cont}}
\providecommand{\Bf}{\mathrm f}
\providecommand{\Bc}{\mathrm c}
\providecommand{\Bcorank}{\operatorname{corank}}
\providecommand{\BPic}{\operatorname{Pic}}
\providecommand{\BCl}{\operatorname{Cl}}
\providecommand{\Br}{\operatorname{Br}}
\providecommand{\BCH}{\operatorname{CH}}
\providecommand{\BRHom}{\operatorname{RHom}}
\providecommand{\BRGamma}{\operatorname{R\Gamma}}

\subsection{Conventions and the audit of the inherited claims}
The symbols in this section distinguish five different theories.  We write
\(H^i_{\BM}(X,\Q(n))\) for rational motivic cohomology in the sense of
Voevodsky (equivalently, the appropriate Adams eigenspace of rational
\(K\)-theory for the regular schemes considered below),
\(H^i_{\et}(X,\Z(n))\) for etale motivic hypercohomology,
\(H^i_{\Bcts}(X,\Q_\ell(n))\) for continuous \(\ell\)-adic cohomology,
\(H^i_{\BW}\) for the specified Weil--etale theory, and
\(H^1_{\Bf}\) for the Bloch--Kato Selmer subgroup.  A subscript \(\Bf\)
in degrees other than one requires the choice of a Selmer complex; it is
not defined just by taking subgroups of ordinary cohomology.  Tate twists
are normalized so that \(\Q(1)\) has weight \(-2\).

For a finite field \(k=\F_q\), arithmetic Frobenius is
\(\phi_q:x\mapsto x^q\), and \(W_k=\langle\phi_q\rangle\cong\Z\).
The generator \(e_q\in H^1(W_k,\Z)\) sends \(\phi_q\) to one.
When using geometric Frobenius on geometric cohomology we use its inverse
consistently; changing that convention changes the sign of \(e_q\), not
the multiplication-by-degree formula below.  In number-ring formulae,
\(\bar X\) means the Artin--Verdier compactification.  Omitting this bar
is a substantive change of theory, even when \(X\) is proper over
\(\Spec\Z\).

The attached transcriptions, rather than scans, are the witnesses used
here.  Thus an asserted correction of a printed formula below is a
mathematical correction of the transcription; it does not constitute a
new paleographic reading.  The reading's R5 is correct about the rational
class in degree three only after the theory is specified.  Its integral
tables omit two nonzero groups of order two.  R6 correctly detects the
failure of unrestricted arithmetic \(\ell\)-adic comparison, but its
positive finite-field formulation needs the additional rational-versus-
numerical-equivalence condition explained below.

\subsection{B.1: the actual finite-field theorem and the missing hypothesis}

\begin{proposition}[Restriction on the arithmetic class]\label{prop:B-e}
\status{Proved}
For \(k_m=\F_{q^m}\), restriction sends
\[
 e_q\longmapsto m e_{q^m}.
\]
For the point, Weil--etale motivic cohomology is
\[
H^i_{\BW}(\Spec\F_q,\Z(n))=
\begin{cases}
\Z,&n=0,\ i=0,1,\\
\Z/(q^{|n|}-1),&n\ne0,\ i=1,\\
0,&\text{otherwise},
\end{cases}\tag{B.1}
\]
where for negative twists the etale convention is
\(\Z(n)=\Q/\Z'(n)[-1]\), with the characteristic-primary part omitted.
After tensoring with \(\Q\), only the two groups with \(n=0\) survive.
\end{proposition}
\begin{proof}
The subgroup \(W_{k_m}\subset W_k\) is generated by \(\phi_q^m\).
The restricted homomorphism has value \(m\) on this generator.  For
\(n=0\), use the length-one resolution of the trivial module over
\(\Z[t,t^{-1}]\), giving \([\Z\xrightarrow{0}\Z]\).
For \(n>0\), the finite-field motivic calculation gives the cyclic group
\(K_{2n-1}(\F_q)=\Z/(q^n-1)\) in degree one and no other group.
For \(n<0\), Frobenius-minus-one is surjective on
\(\Q/\Z'(n)\), and its kernel has order \(q^{|n|}-1\).
The coefficient complex is shifted by minus one.  Weil and etale
cohomology agree for torsion complexes, so the asserted negative-twist
calculation follows.  These comparisons are the torsion and rational
parts of Geisser's Theorem~7.1~\cite{BGeisser2004}.
\end{proof}

The expression \(\Z/(1-q^n)\) in 15:173 is therefore an ordinary cyclic
group only for nonnegative \(n\).  For \(n<0\), its meaningful version is
\(\Z[1/p]/(1-q^n)\), canonically isomorphic to
\(\Z/(q^{|n|}-1)\).  Similarly, the table of continuous cohomology of
\(\widehat\Z\) with coefficients \(\Z_\ell(n)\) presupposes
\(\ell\ne p\); the characteristic-primary coefficient theory is not
the same cyclotomic etale sheaf.

\begin{theorem}[The precise comparison gate]\label{thm:B-Geisser-gate}
\status{Proved}
Let \(X/\F_q\) be smooth and projective of dimension \(d\).  In the
notation of Geisser's Conjectures~8.1--8.3, let \(K(X,n)\) denote the
integral Weil--etale to continuous-cohomology comparison, \(L(X,n)\)
finite generation, and \(T(X,n)\) the conjunction of the Tate cycle-class
conjecture in codimension \(n\), semisimplicity at eigenvalue one, and
equality of rational and numerical equivalence in that codimension.
Then
\[
 K(X,n)+K(X,d-n)\ \Longrightarrow\ L(X,n)
 \ \Longrightarrow\ K(X,n)\ \Longrightarrow\ T(X,n).
 \tag{B.2}
\]
Assuming \(T(Y,r)\) for every smooth projective variety over finite fields
and every \(r\), all the comparisons \(K(X,n)\) and finite-generation
statements \(L(X,n)\) follow.  Moreover, independently of these
conjectures,
\[
 H^i_{\BW}(X,\Q(n))\simeq
 H^i_{\BM}(X,\Q(n))\oplus H^{i-1}_{\BM}(X,\Q(n)).
 \tag{B.3}
\]
Cup product with \(e_q\) acts on these summands by
\( (a,b)\mapsto(0,a)\).
\end{theorem}
\begin{proof}
Equation (B.2) is precisely Geisser's Theorem~8.4, with the three
clauses of his Conjecture~8.3 retained; its last clause must not be
dropped.  Equation (B.3) and the cup-product matrix are Theorem~7.1(c),
deduced there from the computation of the derived pushforward between
the Weil and etale topoi.  Rational etale motivic cohomology agrees with
rational motivic cohomology here.  Thus these are statements about the
actual specified complexes, rather than dimension counts for numerical
cycles~\cite[Theorems~7.1 and~8.4]{BGeisser2004}.
\end{proof}

This identifies an exact defect in the proposed B.1 premise.  The Tate
conjecture and semisimplicity alone supply the expected
\emph{continuous}\(\ell\)-adic groups; they do not, by the cited theorem,
supply the missing equality of rational and numerical equivalence.
Conversely the requested integral comparison already implies that
equality.  No counterexample to that equality is asserted here.  The
unresolved implication must therefore remain part of B.1's status.

For clarity, let \(V^r=H^r_{\et}(\bar X,\Q_\ell(n))\), with
\(\ell\ne p\).  Hochschild--Serre gives
\[
0\longrightarrow V^{i-1}_{\phi_q}
\longrightarrow H^i_{\Bcts}(X,\Q_\ell(n))
\longrightarrow (V^i)^{\phi_q}\longrightarrow0.\tag{B.4}
\]
Purity puts a possible eigenvalue one only in \(r=2n\).
Consequently the groups can be nonzero only in degrees \(2n,2n+1\).
Semisimplicity at one makes the natural invariants-to-coinvariants map
an isomorphism.  Under the Tate conjecture its source is the image of
codimension-\(n\) cycles.  If rational and numerical equivalence agree,
(B.2)--(B.3) give exactly the manuscript's cycle spaces in these two
degrees and the isomorphism by cup product with \(e_q\).  The odd
identification requires the chosen normalization of the arithmetic
generator; this is the valid content of ``pas canoniquement.''

\begin{proposition}[A contradiction already at the point]\label{prop:B-continuity}
\status{Refuted; corrected statement proved}
There is no cohomology theory which simultaneously has the finite-field
Weil--etale value \(H^1(\F_q,\Q)=\Q e_q\), the restriction in
Proposition~\ref{prop:B-e}, continuity under the filtered union of finite
extensions, and vanishing of higher cohomology over \(\overline{\F}_q\).
\end{proposition}
\begin{proof}
All transition maps in the directed system of rational degree-one
groups are multiplication by a positive integer, hence isomorphisms.
Its colimit is a one-dimensional rational space, not zero.  Integrally
the same colimit of copies of \(\Z\), with maps multiplication by the
extension degrees, is \(\Q\).  The contradiction uses no conjecture.
Ordinary etale cohomology of the geometric point instead has no higher
\(\Q_\ell\)-cohomology: it is a different theory and does not repair this
incompatible list of requirements.
\end{proof}

The unqualified demand that actual Weil--etale cohomology satisfy
\emph{every} assertion of 12:60--73 is thus refuted.  The two assertions
must be separated in any replacement.  Nor can the degree-one group be
identified with \(\operatorname{Ext}^1\) in the semisimple category of
pure numerical motives: that extension group vanishes for the unit,
whereas \(H^1_{\BW}(\F_q,\Q)=\Q\).  A theory of extensions must specify
a larger, nonsemisimple arithmetic category or a different derived
construction.

\paragraph{What is unconditionally identified, and what uniqueness can mean.}
For curves all twists satisfy the comparison: this is Geisser's
Theorem~9.5 (dimension at most one belongs to its stated range).
Theorem~9.3 gives the codimension-one comparison whenever the divisor
Tate conjecture holds, in particular for abelian varieties.  For a smooth
projective geometrically connected curve \(C/\F_q\), the two relevant
integral rows can also be checked without appealing to the general
conjectural mechanism:
\[
\begin{array}{c|rrrr}
 &i=0&i=1&i=2&i=3\\ \hline
H^i_{\BW}(C,\Z)&\Z&\Z&\BPic^0(C)^D&(\F_q^\times)^D\\
H^i_{\BW}(C,\Z(1))&0&\F_q^\times&\BPic(C)&\Z .
\end{array}\tag{B.8}
\]
Here \(\Z(1)=\Gm[-1]\), and all unlisted degrees vanish.  This is
Lichtenbaum's Theorem~6.1 after that shift.  The rational rows are
therefore \(\Q\) in degrees zero and one for twist zero and in degrees
two and three for twist one~\cite[Theorem~6.1]{BLichtenbaum2005}.

\begin{theorem}[All twists for motives of abelian type, using the current theorem]
\label{thm:B-abelian-type}
\status{Proved from the cited current theorem}
Let \(X/\F_q\) be smooth projective, and suppose that its rational
Chow motive belongs to the thick rigid category generated by abelian
varieties and Artin motives.  Then, for every \(i,n\) and every prime
\(\ell\),
\[
 H^i_{\BW}(X,\Z(n))\otimes\Z_\ell
 \xrightarrow{\sim}H^i_{\Bcts}(X,\Z_\ell(n)),
\]
and all the integral Weil--etale groups are finitely generated.
For \(\ell=p\), the right side is the logarithmic Hodge--Witt
comparison, rather than a cyclotomic etale sheaf in characteristic
\(p\).  In particular,
\[
 H^i_{\BW}(X,\Q(n))=
 \begin{cases}
 CH^n(X)_\Q,&i=2n,\ 2n+1,\\
 0,&\text{otherwise},
 \end{cases}
\]
where cup product by the normalized \(e_q\) identifies the two
copies.  Thus the previously known divisor and curve cases extend to
all twists on all abelian varieties, and on all varieties of abelian
type.  This statement does not assert that every Chow motive is of
abelian type.
\end{theorem}
\begin{proof}
The external input dated 2026 is van Hoften--Sempliner's
Theorem~I, the Hodge conjecture for CM abelian varieties.  Milne's
Theorem~7.1 then proves the Tate conjecture for every abelian variety
over every finite field~\cite[Theorem~I and \S1.1.3]{AvanHoftenSempliner2026}
\cite[Theorem~7.1]{AMilne1999}.  We use that theorem as an opened
literature result, not as a proof established in this research.

Tate classes pass to direct summands by applying the algebraic
projector to a representing cycle.  They pass to Tate twists and
finite sums, while a tensor product of abelian motives is again the
motive of an abelian variety, up to the finite-extension and twist
operations already allowed.  For a variety over a finite extension,
Galois averaging of an algebraic cycle, followed by pushforward,
descends the required invariant cycle with rational coefficients.
Consequently the Tate statement passes to the whole stated category.
Its Frobenius realizations are semisimple: this holds on \(H^1\)
of an abelian variety, passes to its exterior algebra and all tensor
constructions, and is preserved by projectors and finite extension.
This is also Kahn's Lemma~1.9.

There is an essential extra step beyond Tate surjectivity.  Kahn's
Theorem~1.4 places these Chow motives in the Kimura finite-dimensional
category, and Proposition~1.5 makes the kernel of the numerical
endomorphism map nilpotent.  His Theorem~1.10 therefore identifies
rational and numerical equivalence here.  Corollary~2.2 gives the
vanishing of positive rational \(K\)-groups.  With these hypotheses
verified, his Theorem~3.6 proves both the \(\ell\ne p\) and
the logarithmic \(p\)-adic comparisons; Corollary~3.8(c) supplies
finite generation.  The rational formula follows either from his
Lemma~3.7 and Corollary~3.8(d), or by substituting rational
motivic vanishing into Geisser's splitting (B.3).
All these are exact statements for the same category, so no general
Beilinson assumption has entered~\cite[Theorem~1.10, Corollary~2.2,
Theorem~3.6, Corollary~3.8]{BKahn2003}.
\end{proof}

\begin{proposition}[The remaining universal assertion is exactly a phantom-motive gate]
\label{prop:B-phantom-gate}
\status{Proved equivalence; the equivalent universal assertion remains unresolved}
Assume Tate surjectivity and semisimplicity at eigenvalue one for
every smooth projective variety over \(\F_q\), in every
codimension.  Then the following universal statements are equivalent:
\begin{enumerate}
\item The integral Weil--etale comparison holds for all varieties,
degrees and twists, and all these integral groups are finitely generated.
\item Rational equivalence equals numerical equivalence on all
rational Chow groups.
\item There is no nonzero rational Chow motive \(P\) whose image
\(\bar P\) in numerical motives is zero.
\item Every rational Chow motive belongs to the category generated
by abelian and Artin motives.
\end{enumerate}
In particular, the missing assertion in B.1 can be expressed without
an unspecified regulator or unspecified comparison map: it is the
vanishing of numerical Chow phantoms.  The Tate and semisimplicity
hypotheses alone have not been shown to imply that vanishing here.
\end{proposition}
\begin{proof}
The precise equivalences, including integral comparison and finite
generation, are Kahn's Theorems~5.1 and~5.2, especially clauses
(iv bis), (iv ter), (v), (vi ter), and (vi quater).  Here is the
categorical obstruction in explicit form.  Strong Tate, which is
the stated combination of cycle surjectivity and semisimplicity,
implies that every numerical motive is of abelian type.  Thus, for
a Chow motive \(M\), there is an abelian-type Chow motive \(N\)
and inverse numerical isomorphisms
\(\bar f:\bar M\to\bar N\),
\(\bar g:\bar N\to\bar M\).
Fullness of the numerical quotient lifts them to \(f,g\).
The element \(u=1_N-fg\) is numerically zero.  Nilpotence on
the abelian-type motive \(N\) gives \(u^r=0\) for some \(r\).
Replace \(f\) by
\[
 f'=(1_N+u+\cdots+u^{r-1})f.
\]
Then \(f'g=1_N\), so \(gf'\) is an idempotent on \(M\).
The rational Chow category splits idempotents, hence
\[
 M\simeq N\oplus P,\qquad \bar P=0.
\]
This constructs the exact leftover motive.  If numerical phantoms
vanish, every \(M\) is abelian-type and
Theorem~\ref{thm:B-abelian-type} applies.  Conversely universal
rational--numerical equality makes the numerical quotient faithful:
if \(\bar P=0\), then \(1_P=0\), so \(P=0\).
The equivalence of that equality with universal comparison is also
Geisser's Theorem~8.4 after the Tate and semisimplicity hypotheses
have been fixed.  Kahn's argument identifies precisely which
additional step the two stated hypotheses do not supply
\cite[Theorems~5.1--5.2 and the final proof on pp.~999--1000]{BKahn2003}.
\end{proof}

There is a useful limited uniqueness statement.  Fix a semisimple
\(\Q\)-linear category of pure motives with simple unit \(\mathbf1\)
and \(\End(\mathbf1)=\Q\), and assume its realization detects precisely
the unit-isotypic part as the Frobenius-one eigenspace.  Any
\(\Q\)-linear cohomological functor on its bounded derived category
with the prescribed continuous-cohomology values on pure objects is,
noncanonically, the functor
\[
 M\longmapsto\Hom(\mathbf1,M)\oplus
 \bigl(\Hom(\mathbf1,M)\otimes L\bigr)[-1],\qquad \dim_\Q L=1.
 \tag{B.9}
\]
Indeed every object is a direct sum of simple objects; comparison kills
all nonunit simples in every degree and prescribes the two groups of
the unit.  Every complex in a semisimple category splits into its
cohomology, so there are no further extension choices for this additive
derived functor.  A chosen \(e\in L\setminus0\) normalizes the
isomorphism.  Without \(e\), multiplication by \(a\in\Q^\times\) on
the odd factor is a nontrivial natural automorphism preserving \(H^0\).
This proves uniqueness only in this specified additive sense, for this
category.  It is not a uniqueness theorem for topoi, integral complexes,
products, six operations, or rational cycle classes.  Identifying (B.9)
with actual motivic Weil--etale cohomology requires the gate in
Theorem~\ref{thm:B-Geisser-gate}; the manuscript's unrestricted list is
already inconsistent by Proposition~\ref{prop:B-continuity}.

\subsection{B.2: Tate twists on number rings and the exact regulator defects}
Let \(K\) have \(r_1\) real and \(r_2\) complex places, let \(S\) be
a finite set of finite places, and put
\[
 X=\Spec\mathcal O_{K,S},\quad T=S\cup S_\ell,\quad
 U=X[1/\ell]=\Spec\mathcal O_{K,T},\quad
 r_n=\begin{cases}r_2&n\text{ even},\\r_1+r_2&n\text{ odd}.
 \end{cases}
\]
All the cohomology in the next table is \emph{full continuous}
cohomology of \(U\).  In degree one it is equivalently the cohomology
of the Galois group of the maximal extension unramified outside
\(T\) and infinity.

There is a scope correction to R6's proposed comparison example.
Its ranks one and zero for \(\Q(0)\) compare two constant-coefficient
theories on the fixed scheme \(\Spec\Z[1/p]\).  If, as B.2
explicitly requests, the scheme is instead replaced by
\(X[1/\ell]\) for each coefficient prime, the degree-one ranks in
twist zero are both one.  A counterexample with exactly this latter
convention is \(X=\Spec\Z[1/p]\) and twist one: (B.11) gives
rank one for \(\ell=p\), but rank two for \(\ell\ne p\).
The extra generator is the newly inverted prime \(\ell\).
Thus full cohomology still fails the requested common rational
comparison, but its verification must use the same base convention.

\begin{theorem}[All full rational etale groups, with the defect left visible]
\label{thm:B-full-tate}
\status{Proved}
For every prime \(\ell\),
\[
\begin{array}{c|ccc}
n&H^0_{\Bcts}(U,\Q_\ell(n))&
\dim_{\Q_\ell}H^1_{\Bcts}(U,\Q_\ell(n))&
\dim_{\Q_\ell}H^2_{\Bcts}(U,\Q_\ell(n))\\ \hline
0&\Q_\ell&r_2+1+\delta_{K,\ell}(0)&\delta_{K,\ell}(0)\\
1&0&r_1+r_2-1+|T|&|T|-1\\
n\ge2&0&r_n&0\\
n<0&0&r_n+\delta_{K,\ell}(n)&\delta_{K,\ell}(n).
\end{array}\tag{B.10}
\]
All other degrees vanish.  For \(n=1\), more precisely,
\[
 H^1_{\Bcts}(U,\Q_\ell(1))=
 \mathcal O_{K,T}^\times\otimes\Q_\ell,\qquad
 H^2_{\Bcts}(U,\Q_\ell(1))=
 \ker\bigl(\Q_\ell^T\xrightarrow{\sum}\Q_\ell\bigr).
 \tag{B.11}
\]
For \(n\ge2\), the Chern-class map identifies degree one with
\(K_{2n-1}(\mathcal O_K)^{(n)}\otimes\Q_\ell\).
The defect \(\delta_{K,\ell}(0)\) is exactly the Leopoldt defect.
For \(n<0\), writing \(m=1-n\), it is exactly the dimension of
the kernel of the \(\ell\)-adic regulator
\[
 K_{2m-1}(\mathcal O_K)^{(m)}\otimes\Q_\ell
 \longrightarrow\bigoplus_{v\mid\ell}K_v .\tag{B.12}
\]
In particular, replacing every defect in (B.10) by zero is equivalent,
at the indicated \((K,\ell,n)\), to Leopoldt for \(n=0\), or to
Schneider's regulator-injectivity conjecture at \(m=1-n\) for \(n<0\).
\end{theorem}
\begin{proof}
For twist one, Kummer theory identifies degree one with units because
the ideal class group is finite.  The Brauer reciprocity sequence
identifies the rational Tate module of \(\Br(U)\) with the kernel of
the sum of the finite local invariants; real places contribute only
groups killed by two.  This proves (B.11).

For positive twists at least two, the rational Chern-class comparison
and the finiteness of the degree-two integral cohomology give the
assertions; the rank is Borel's \(r_n\).  For odd \(\ell\) a precise
modern reference is Nickel's Theorem~3.2; at two use Rognes--Weibel's
Theorems~0.4 and~0.6, after rationalization.  Their real-place errors
are finite two-groups, hence do not alter this rational assertion.
Localization away from \(\ell\) does not change either group when
\(n\ne1\), since the additional residue-field terms have twist
\(n-1\ne0\) and are rationally zero~\cite[Theorem~3.2]{BNickel2022}
\cite[Theorems~0.4 and~0.6]{BRognesWeibel2000}.

For \(n\le0\), Poitou--Tate duality identifies
\[
 H^2_{\Bcts}(U,\Q_\ell(n))^\vee=
 \ker\left(H^1_{\Bcts}(U,\Q_\ell(1-n))
 \longrightarrow\bigoplus_{v\in T}H^1(K_v,\Q_\ell(1-n))\right).
 \tag{B.13}
\]
No local degree-two term intervenes: its dual is
\(H^0(K_v,\Q_\ell(1-n))=0\).  If \(n=0\), the valuations first
restrict this kernel to \(\mathcal O_K^\times\otimes\Q_\ell\), and
the remaining map is the vector of local \(\ell\)-adic logarithms.
This is exactly the definition of the Leopoldt defect.  If \(n<0\),
the positive-twist comparison and the inverse Bloch--Kato exponential
identify the map with (B.12).  These descriptions also show independence
of the extra places in \(S\).  The global Euler characteristic is
\(-r_n\), and \(H^0\) is one-dimensional only in twist zero; hence
the degree-one dimensions follow.  Rational cohomological dimension is
two, including at \(\ell=2\), since the additional real-place
cohomology is two-torsion.  Finally, Nickel's Conjecture~3.5,
Remarks~3.7--3.10 and Lemma~3.11 identify precisely this strict
localization kernel with Schneider's conjecture; no nonvanishing of
an unspecified regulator has been assumed~\cite[\S3.6]{BNickel2022}.
\end{proof}

\paragraph{The local conditions must depend on the arithmetic open.}
For \(L/\Q_\ell\) finite, \(d_L=[L:\Q_\ell]\), the dimensions of
the Bloch--Kato local subgroups are
\[
\begin{array}{c|rrrr}
 n&\dim H^1_e&\dim H^1_f&\dim H^1_g&\dim H^1\\ \hline
n<0&0&0&0&d_L\\
0&0&1&1&d_L+1\\
1&d_L&d_L&d_L+1&d_L+1\\
n>1&d_L&d_L&d_L&d_L .
\end{array}\tag{B.14}
\]
The twist-zero finite line consists of unramified characters; the
twist-one finite subgroup is \(\mathcal O_L^\times\otimes\Q_\ell\),
while its geometric subgroup is all of
\(L^\times\widehat\otimes\Q_\ell\).  For a finite place
\(v\nmid\ell\), the full \(H^1(K_v,\Q_\ell(n))\) has dimension
one exactly for \(n=0,1\), and zero otherwise; its unramified subgroup
has dimension one exactly for \(n=0\).  The first table is
Bloch--Kato's Example~3.9, with their Proposition~3.8 providing the
orthogonality statements.  The second follows directly from inertia
cohomology and the eigenvalue of Frobenius~\cite[Proposition~3.8,
Example~3.9]{BBlochKato1990}.

Define \(H^1_{f,X}(K,V)\) using \(H^1_f(K_v,V)\) at finite
\(v\notin S\) and \(H^1_g(K_v,V)\) at \(v\in S\); this is
exactly the convention of Bloch--Kato's Definition~5.1.  Away from
\(\ell\), the latter condition imposes no restriction for the Tate
representations at issue.  Imposing \(f\) also at a removed place
above \(\ell\) would incorrectly discard that place's unit valuation.
Imposing no condition on twist zero there would retain precisely the
ramified characters responsible for R6.

\begin{proposition}[An unconditional rational structure in degree one]
\label{prop:B-Hf-tate}
\status{Proved}
For every integer \(n\), all primes \(\ell\), and the above
geometric local conditions, the natural realization maps give
\[
 H^1_{f,X}(K,\Q_\ell(n))\simeq V_X(n)\otimes\Q_\ell,
\quad
 V_X(n)=
 \begin{cases}
 \mathcal O_{K,S}^\times\otimes\Q,&n=1,\\
 K_{2n-1}(\mathcal O_K)^{(n)}\otimes\Q,&n\ge2,\\
 0,&n\le0.
 \end{cases}\tag{B.15}
\]
In particular the degree-one assertion for negative twists does not
require Schneider's conjecture.  That conjecture concerns the full
groups in (B.10), not the group in (B.15).
\end{proposition}
\begin{proof}
For \(n=1\), the Kummer map in (B.11) and (B.14) keep exactly
the \(S\)-units.  For \(n\ge2\), all local finite conditions are
already satisfied, so Theorem~\ref{thm:B-full-tate} applies.  For
\(n=0\), classes satisfying these conditions are continuous
\(\Q_\ell\)-valued characters unramified at every finite place.
Global class field theory factors them through a finite ideal class
group (and finite real-place factors); hence they vanish.  For \(n<0\),
both local \(f\) and \(g\) are zero.  The strict global kernel is
dual, by (B.13) with the twists interchanged, to the degree-two group
in positive twist \(1-n\ge2\), which is zero by the positive-twist
theorem.  This proves the assertion without setting
\(\delta_{K,\ell}(n)\) to zero.
\end{proof}

\paragraph{Degrees other than one and the scope of a corrected requirement 4).}
On the proper number ring there is a precise higher-degree correction.
Take the Bloch--Kato Selmer complex.  At a place over \(\ell\),
the local finite complex for the crystalline Tate representation is
\[
 C_f(K_v,V)=
 [D_{\rm cris}(V)\xrightarrow{\ ((1-\varphi),\mathrm{can})\ }
 D_{\rm cris}(V)\oplus t_V],\qquad
 t_V=D_{\rm dR}(V)/\operatorname{Fil}^0,
\]
in degrees zero and one.  Away from \(\ell\), use
\([V^{I_v}\xrightarrow{1-\operatorname{Frob}_v}V^{I_v}]\).
These complexes have cohomology \(H^0(K_v,V)\) and
\(H^1_f(K_v,V)\), and no other cohomology.  The global complex is
the mapping fibre of the global cochain complex plus these local
complexes mapping to the direct sum of the full local complexes.
Fontaine's \S\S4.4--4.5 give these constructions and their duality;
the real-place Tate errors are two-torsion and disappear in this
rational calculation~\cite[pp.~214--216]{BFontaine1992}.

Here is the duality argument in the degrees actually used.  Local
Tate duality and Bloch--Kato's orthogonality of finite subspaces
give an isomorphism of the quotient cone with
\[
 \operatorname{Cone}\bigl(C_f(K_v,V)\longrightarrow
   R\Gamma(K_v,V)\bigr)
 \simeq R\operatorname{Hom}_{\Q_\ell}
       (C_f(K_v,V^\vee(1)),\Q_\ell)[-2].
\]
Indeed its degree-one quotient is dual to the complementary finite
subspace, and its degree-two local invariant is dual to
\(H^0(K_v,V^\vee(1))\); all other groups vanish.
To keep the issue of choices explicit, set
\[
\begin{split}
 C(V)&=\operatorname{coker}\left(H^1(U,V)\longrightarrow
       \bigoplus_{v\in T}H^1(K_v,V)/H^1_f(K_v,V)\right),\\
 K(V)&=\ker\left(H^2(U,V)\longrightarrow
       \bigoplus_{v\in T}H^2(K_v,V)\right).
\end{split}
\]
The Selmer cone and the Poitou--Tate sequence, respectively, give
\[
\begin{gathered}
0\longrightarrow C(V)\longrightarrow H_f^2(K,V)
 \longrightarrow K(V)\longrightarrow0,\\
0\longrightarrow C(V)\longrightarrow
 H_f^1(K,V^\vee(1))^\vee
 \longrightarrow K(V)\longrightarrow0.
\end{gathered}
\]
These sequences split as vector spaces.  Thus the degree-two
isomorphism below exists, but is not asserted to be canonical
from the finite subgroups alone.  The extreme-degree
identifications follow directly from global invariants and the
local invariant sum.  We obtain
\[
 H_f^0=H^0(K,V),\quad H_f^1=H_f^1(K,V),\quad
 H_f^2\simeq H_f^1(K,V^\vee(1))^\vee,\quad
 H_f^3\simeq H^0(K,V^\vee(1))^\vee.\tag{B.16}
\]
For Tate motives, (B.15)--(B.16) supply a common rational space in
every degree, including the dual unit and dual Borel-regulator spaces;
the higher comparison isomorphisms use the specified splittings.
Their graded groups agree with Flach--Morin's rational compactified
complex (B.6), now for \(\mathcal O_K\).  This also proves why a comparison with
\emph{ordinary} motivic cohomology in every degree would be false:
\(H^3_{\BM}(\Spec\Z,\Q(1))=0\), whereas (B.16) gives \(\Q_\ell\).

For an arithmetic open, merely specifying the subgroups in degree one
does not specify the higher Selmer complex.  In particular, at a
removed \(\ell\)-adic place in twist one one must retain the local
degree-two invariant as well as \(H^1_g=H^1\) to reproduce the
localization triangle.  The rational absolute complex with localization
has, in addition to the degree-one \(S\)-unit group,
\[
 H^2(X,\Q(1))=\ker(\Q^S\xrightarrow{\sum}\Q),\qquad
 H^3(X,\Q(1))=\operatorname{coker}(\Q^S\xrightarrow{\sum}\Q).
 \tag{B.17}
\]
For \(n\ne1\) it has the same groups as the compactified proper
number ring.  This follows by applying the finite-field computation
to the Gysin boundary in twist \(n-1\).  It is a construction of
a specified rational absolute theory; it is not an assertion that
every expression denoted \(H^i_f\) in the literature has these values.

\begin{corollary}[The common rational spaces in every Tate degree]
\label{cor:B-Hf-all-degrees}
\status{Proved for the specified absolute and Selmer complexes}
For \(X=\Spec\mathcal O_{K,S}\), the following table lists all
possibly nonzero common \(\Q\)-spaces \(H^i_{\rm abs}(X,\Q(n))\).
The dual in the table is the \(\Q\)-linear dual.
\[
\begin{array}{c|c|l}
 n&i&H^i_{\rm abs}(X,\Q(n))\\ \hline
 0&0&\Q\\
 0&2&(\mathcal O_K^\times\otimes\Q)^\vee\\
 1&1&\mathcal O_{K,S}^\times\otimes\Q\\
 1&2&\ker(\Q^S\xrightarrow{\sum}\Q)\\
 1&3&\operatorname{coker}(\Q^S\xrightarrow{\sum}\Q)\\
 n\ge2&1&K_{2n-1}(\mathcal O_K)^{(n)}\otimes\Q\\
 n<0&2&(K_{1-2n}(\mathcal O_K)^{(1-n)}\otimes\Q)^\vee .
\end{array}
\tag{B.17a}
\]
Every unlisted group is zero.  For every prime \(\ell\), including
two, these spaces tensor to the cohomology of the specified Selmer
complex with the geometric open-place conditions.  Neither Leopoldt
nor Schneider is assumed in this assertion.
\end{corollary}
\begin{proof}
For the proper number ring, substitute (B.15) into (B.16).  For the
open, the local change from the finite to the geometric condition
is zero in every twist except one.  In twist one it adds the local
valuation and invariant in degrees one and two.  The map
\(\mathcal O_{K,S}^\times\otimes\Q\to\Q^S\) is surjective:
its integral cokernel embeds in the finite ideal class group.
The map between the remaining invariant lines is the sum map by
Brauer reciprocity.  This proves exactly (B.17), including
\(S=\varnothing\), when its cokernel is \(\Q\).  Tensoring the
construction with \(\Q_\ell\) gives the local complexes just
described.  The negative-twist assertion uses the dual of the proved
positive-twist group, not injectivity of the regulator (B.12).
\end{proof}

The precise degree-one replacement for the manuscript's requirement
4) is the Bloch--Kato realization conjecture: for a geometric motive
\(M\) over \(K\), define the motivic extension space over an
arithmetic open by the image from a regular model (and use
homologically trivial cycles in the middle case).  Its etale regulator
must induce
\[
 H^1_{\rm mot,f,X}(K,M)\otimes\Q_\ell
 \xrightarrow{\sim}H^1_{f,X}(K,M_\ell)\quad\text{for every }\ell.
 \tag{B.18}
\]
For motives of the form \(h^m(Y)(r)\), this is exactly the
model-image and realization assertion in Bloch--Kato's
Conjecture~5.3(i)--(iii), with the local convention of Definition~5.1.
In Fontaine's direct precursor to his joint work with Perrin-Riou,
\S6.5 expresses the proper-base comparison as \emph{f-admissibility}:
the spaces \(H_f^0(K,M)\) and \(H_f^1(K,M)\) are finite-dimensional,
and their realization maps \(c^0_\ell,c^1_\ell\) are isomorphisms
after tensoring with \(\Q_\ell\), for every \(\ell\).
His \S6.6 conjectures this for the motivic category, and \S6.7
spells it out for the geometrically defined motivic spaces
\cite[pp.~221--223]{BFontaine1992}.  This is a directly opened
primary-source formulation; the later full Fontaine--Perrin-Riou
article was not obtained in this run.
A statement in all degrees additionally needs a chosen arithmetic
motivic category and the dual completion (B.16)--(B.17), as well as
compatibility with localization and duality.  Over finite fields the
corresponding absolute complex is the Weil complex of B.1; its cycle
comparison retains all three clauses of Theorem~\ref{thm:B-Geisser-gate}.
Bloch--Kato's degree-one conjecture alone does not imply those extra
finite-field clauses.  Thus the requested single implication from
an unspecified ``finite condition'' to every requirement of B.1
cannot be asserted as a theorem.

\paragraph{Scholbach's construction.}
Scholbach's Theorem~6.8 identifies an intermediate-extension motivic
group with \(H^1_{f,\backslash\mathrm{crys}}\), under the stated
axioms for mixed motives and their weights.  His Theorems~6.11 and~6.13
identify it with the model-image definition.  These numbers were
checked against the opened arXiv version.  Remark~6.2 explicitly
excludes the crystalline condition at places over \(\ell\) from the
paper's treatment.  The derived \(\Hom\) definition and assumed six
operations support analogues of requirements 1)--3) in its stated
motivic category; the existence and properties of the motivic
\(t\)-structure are among the assumptions.  They do not prove (B.18)
with the missing local conditions restored.  Moreover ordinary
motivic cohomology in that construction has
\(H^1_{\BM}(\F_q,\Q)=0\) and
\(H^3_{\BM}(\Spec\Z,\Q(1))=0\), so it cannot itself be the
full absolute theory with the odd Weil--etale classes.  This is an
explicit distinction, rather than an unidentified comparison
map~\cite[Remark~6.2 and Theorems~6.8, 6.11, 6.13]{BScholbach2015}.

The other supplied sources do not remove this distinction.  In the
published version of Scholbach's \emph{Special L-values}, Axiom~1.6
states the mixed-motive assumptions, and Theorem~4.5 identifies the
finite-field perfectness assertion with rational=numerical equivalence
plus Beilinson--Parshin vanishing~\cite{BScholbachSpecial2017}.
Adrien Morin's supplied constructible-sheaf paper has a different
coefficient scope: Theorem~2.8 treats twists \(n\le0\), under
its stated perfectness hypothesis (or using Weil--eh cohomology and
resolution), and Theorem~8.3 specifies when a motivic application is
of that constructible kind.  These results do not supply the omitted
positive-twist or arbitrary pure-motive realization
\cite{BAdrienMorin2024}.

\subsection{B.3: abelian schemes, local coranks, and the certified counterexample}
Let \(A\) be an abelian scheme on
\(X=\Spec\mathcal O_K[1/N]\), with \(\ell\mid N\), and write
\(A^t\) for its dual.  Let
\(\mathcal S_\ell(A^t)=\varprojlim_a\operatorname{Sel}_{\ell^a}(A^t/K)\)
be the pro-\(\ell\) Selmer group, and let \(D\) denote Pontryagin
duality to \(\Q_\ell/\Z_\ell\).  At real places the local degree-zero
group used in this duality is the component group, as in local Tate
duality.

\begin{theorem}[The exact support sequence and corank]\label{thm:B-Sha-corank}
\status{Proved}
There is an exact sequence
\[
\begin{split}
0\longrightarrow\Sha(A/K)[\ell^\infty]
\longrightarrow H^1_{\et}(X,A)[\ell^\infty]
\longrightarrow\bigoplus_{v\mid N\ \mathrm{or}\ v\mid\infty}
 H^1(K_v,A)[\ell^\infty]
\longrightarrow\mathcal S_\ell(A^t)^D .
\end{split}\tag{B.19}
\]
There is no assertion of surjectivity onto the last term.
Put \(d=[K:\Q]\dim A\),
\(s=\Bcorank_{\Z_\ell}\Sha(A/K)[\ell^\infty]\), and let
\[
 t=\operatorname{rank}_{\Z_\ell}\operatorname{im}
 \left(\mathcal S_\ell(A^t)\longrightarrow
 \prod_{v\mid\ell}A^t(K_v)^{\wedge}_\ell\right).
\]
Then
\[
 \Bcorank_{\Z_\ell}H^1_{\et}(X,A)[\ell^\infty]=d+s-t.
 \tag{B.20}
\]
If \(\Sha(A/K)[\ell^\infty]\) is finite, then \(t\) is the
dimension of the \(\ell\)-adic closure of the Mordell--Weil group
\(A^t(K)\) inside its local points.  It must not in general be
replaced by the abstract rank \(r=\operatorname{rank}A(K)\).
\end{theorem}
\begin{proof}
The generic-fibre map embeds \(H^1(X,A)\) into \(H^1(K,A)\).
Its image consists of classes locally trivial at the finite places
outside \(N\): over those henselian integer rings proper smooth
extension and Lang's theorem give \(H^1(\mathcal O_v,A)=0\).
The converse follows by extending a locally trivial torsor across the
missing codimension-one points and descent on the Dedekind base.
The unconditional Cassels--Tate sequence of Gonz\'alez-Avil\'es and
Tan, Main Theorem, is
\[
0\longrightarrow\Sha(A)[\ell^\infty]\longrightarrow H^1(K,A)[\ell^\infty]
\longrightarrow\bigoplus_v H^1(K_v,A)[\ell^\infty]
\longrightarrow\mathcal S_\ell(A^t)^D\longrightarrow0.
\]
Taking the inverse image of the direct sum supported on \(N\) and
infinity proves (B.19).  This source is needed because replacing
the last term by \(A^t(K)^{\wedge,D}_\ell\) without a finiteness
assumption is unjustified~\cite[Main Theorem and equation~(6)]{BGonzalezTan2007}.

Local Tate duality identifies the dual of each local \(H^1\) with
\(A^t(K_v)^{\wedge}_\ell\).  Its rank is
\([K_v:\Q_\ell]\dim A\) for \(v\mid\ell\), and zero at all
other finite or real places.  The dual of the last map of (B.19)
is precisely the local-point map defining \(t\), so its image has
corank \(t\).  Additivity of coranks in the resulting short exact
sequences gives (B.20).  Finally the Selmer Kummer sequences give
\[
0\longrightarrow A^t(K)^{\wedge}_\ell
\longrightarrow\mathcal S_\ell(A^t)
\longrightarrow T_\ell\Sha(A^t/K)\longrightarrow0.
\]
The dual abelian varieties are isogenous via a polarization, so their
Sha groups have the same corank.  If that corank is zero, these
cofinitely generated groups are finite and their Tate modules vanish.
This proves the last assertion.
\end{proof}

The distinction between \(t\) and \(r\) is already effective over \(\Q\).
The second curve in the attached computation is
\[
 E':y^2+y=x^3+x^2-2x,
 \qquad \Delta(E')=389.
\]
The certified output is \(\operatorname{ellrank}(E')=[2,2,0,
[[-2,0],[4,8]]]\), with trivial torsion.  Its two-Selmer group
therefore has dimension two and \(\Sha(E')[2]=0\).  On
\(\Spec\Z[1/(2\cdot389)]\), we have \(d=1,s=0,r=2\), while
\(t=1\).  To see the last equality without a numerical logarithm,
a non-torsion rational point cannot become torsion in \(E'(\Q_2)\),
and its multiple in the formal-group neighborhood has nonzero
\(2\)-adic logarithm; the closure it generates has rank one.
Consequently (B.20) gives corank zero.  The putative formula
\(d-r+s\) would give minus one and is false.  The same
\texttt{x011.gp} file records all these data.

\begin{example}[The counterexample in R6, without a BSD assumption]\label{ex:B-X011}
\status{Refuted; corrected statement proved}
Take
\[
 E=X_0(11):\quad y^2+y=x^3-x^2-10x-20.
\]
Its discriminant is \(-11^5\), its conductor is \(11\), and it has
good reduction over \(\Z[1/11]\).  The reproducible algebraic
two-descent in \texttt{part\_b/code/x011.gp} gives
\[
 \operatorname{ellrank}(E)=[0,0,0,[]],\qquad
 E(\Q)_{\rm tors}\cong\Z/5,
\]
with torsion generator \((5,5)\).  Thus its rank is certified to be
zero; this uses no analytic-rank computation or BSD.  PARI's official
description of \texttt{ellrank} specifies that the lower and upper
bounds and the two-descent data are unconditional.  The model of
\(X_0(11)\) is also derived explicitly in Elkies's October~13, 2021
lecture notes~\cite{BElkies2021,BPari2026}.

For any prime \(\ell\), let \(X=\Spec\Z[1/(11\ell)]\).
The Kummer sequence gives
\[
0\longrightarrow E(\Q)\otimes\Q_\ell
\longrightarrow H^1_{\Bcts}(X,V_\ell E)
\longrightarrow V_\ell H^1_{\et}(X,E)\longrightarrow0.
\]
The first term is zero.  The global Euler characteristic of
\(V_\ell E\), whose complex-conjugation eigenspaces both have
dimension one, is \(-1\).  Its \(H^0\) is zero, so
\(\dim H^1=1+\dim H^2\ge1\).  It follows that
\(H^1_{\et}(X,E)[\ell^\infty]\) is infinite for every \(\ell\).
For \(\ell=2\) the zero two-Selmer calculation additionally gives
\(\Sha(E)[2]=0\), hence \(\Sha(E)[2^\infty]=0\); then (B.20)
gives exact corank one.  All hypotheses of the boxed assertion on
12:75 hold, including invertibility of \(\ell\) on \(X\), and its
conclusion fails.
\end{example}

\paragraph{The corrected boxed assertion and its precise remaining conjecture.}
Set
\[
 H^1_{\rm loc\,0}(X,A)[\ell^\infty]=
 \ker\left(H^1(X,A)[\ell^\infty]\longrightarrow
 \bigoplus_{v\mid\ell}H^1(K_v,A)[\ell^\infty]\right).
\]
The local groups at the remaining finitely many places in (B.19) are
finite, so this kernel is finite if and only if
\(\Sha(A/K)[\ell^\infty]\) is finite.  The corrected statement is
therefore a Sha-finiteness conjecture, rather than a theorem obtained
by putting \(\ell\) in the denominator.  At the rational Tate-module
level, local Kummer theory identifies the finite condition with
the image of \(A(K_v)\widehat\otimes\Q_\ell\) (Bloch--Kato,
Example~3.11).  The global Selmer Kummer sequence becomes
\[
0\longrightarrow A(K)\otimes\Q_\ell
\longrightarrow H_f^1(K,V_\ell A)
\longrightarrow V_\ell\Sha(A/K)\longrightarrow0.\tag{B.21}
\]
Thus the rational-points assertion, with the restored finite local
conditions, is equivalent exactly to the finite boxed kernel.
This is the content of Bloch--Kato's Proposition~5.4, here derived
from Kummer theory.  The motive written \(T_{\rm mot}(A)\) on
12:74 has realization \(V_\ell A\) and weight \(-1\), so the
transcription's weight \(+1\) must also be corrected.

\begin{proposition}[An arithmetic relative-curve counterexample to 12:79]
\label{prop:B-Brauer-arithmetic}
\status{Refuted; corrected statement proved}
Let \(S=\Spec\Z[1/22]\) and let \(f:X\to S\) be the smooth
proper genus-one curve given by the above model of \(X_0(11)\), with
its zero section.  Then
\[
 \left(\Br(X)/f^*\Br(S)\right)[2^\infty]
 \simeq H^1_{\et}(S,E)[2^\infty]
\]
has \(\Z_2\)-corank one.  Thus surjectivity modulo finite groups
of \(\Br(S)\to\Br(X)\) is false for arithmetic bases, even for
a smooth proper relative curve with a section.
\end{proposition}
\begin{proof}
The zero section splits the relative degree map, giving
\(R^1f_*\Gm=\BPic_{X/S}\simeq E\times\Z\) as etale sheaves.
Proper base change identifies \(R^2f_*\mu_{2^a}\) with the
constant sheaf \(\Z/2^a\); the degree-one section class maps to its
generator.  The relative Kummer sequence therefore shows that
\((R^2f_*\Gm)[2^a]=0\) for every \(a\).  The Leray
spectral sequence, and the splitting on pullback supplied by the
section, now give on the torsion subgroup the usual exact edge
sequence
\[
 0\longrightarrow\Br(S)\longrightarrow\Br(X)
 \longrightarrow H^1(S,\BPic_{X/S})
 \longrightarrow H^3(S,\Gm),
\]
as far as its \(2\)-primary component is concerned.  Here
\(H^1(S,\Z)=0\), since a continuous homomorphism from the profinite
etale fundamental group to the discrete torsion-free group \(\Z\)
is zero.  Moreover \(H^3(S,\Gm)=0\): Milne's
Proposition~II.2.1 gives its description as the cokernel of the
sum of local invariants at the omitted places, and the summand at
the omitted finite place \(2\) already surjects onto \(\Q/\Z\).
This is stated explicitly in Remark~II.2.2(a)~\cite{AMilneADT}.
All groups in \(H^1(S,E)\) are torsion, and \(\Br(X)\) is
torsion because \(X\) is regular.  Taking primary components of
the resulting exact sequence is therefore exact.  The preceding
certified two-descent and (B.20) give the asserted corank one.
\end{proof}

\paragraph{The finite-field Hochschild--Serre correction.}
Let now \(S/\F_q\) be smooth and geometrically connected,
\(A/S\) an abelian scheme, \(\ell\ne p\),
\(\Gamma=\Gal(\overline{\F}_q/\F_q)\), and
\(C=A(\bar S)\).  The actual five-term sequence is
\[
0\longrightarrow H^1(\Gamma,C)\longrightarrow H^1(S,A)
\longrightarrow H^1(\bar S,\bar A)^\Gamma
\xrightarrow{d_2}H^2(\Gamma,C).\tag{B.29}
\]
The last term deleted on 12:76--77 cannot be discarded solely
because \(\Gamma\) has torsion cohomological dimension one: \(C\)
is not a torsion module.  For a lattice \(M\) with continuous
\(\Gamma\)-action, the exact calculation is
\[
0\longrightarrow M^\Gamma\otimes\Q/\Z
\longrightarrow(M\otimes\Q/\Z)^\Gamma
\longrightarrow H^1(\Gamma,M)\longrightarrow0,
\qquad
H^2(\Gamma,M)=(M\otimes\Q/\Z)_\Gamma.\tag{B.30}
\]
Indeed \(M\otimes\Q\) has zero higher continuous cohomology:
every cochain factors through a finite quotient and averaging is
exact over \(\Q\).  The rational-to-torsion exact sequence then
proves (B.30).  In particular \(H^1(\Gamma,M)\) is finite,
whereas \(M=\Z\) gives \(H^2(\Gamma,M)=\Q/\Z\).
The equality with \((M\otimes\Q/\Z)^\Gamma\) asserted for
\(H^1(\Gamma,M)\) on 12:76 misses the first term of (B.30).

Modulo the constant trace, Lang--Neron gives the finite-rank group
\(M\) on the page.  Lang's theorem and torsion cohomological
dimension one dispose of the trace's higher cohomology, so
\(H^1(\Gamma,C)\) is finite; they do not dispose of \(d_2\).
Write
\[
 W=H^1_{\Bcts}(\bar S,V_\ell A),\qquad
 M_\ell=C^{\wedge}_\ell\otimes_{\Z_\ell}\Q_\ell
\]
where \(C^{\wedge}_\ell=\varprojlim_a C/\ell^a C\), so the
divisible trace is removed.  Geometric Kummer
gives
\[
 0\longrightarrow M_\ell\longrightarrow W
 \longrightarrow V_\ell H^1(\bar S,\bar A)
 \longrightarrow0.\tag{B.31}
\]
Its invariant sequence includes the boundary
\((V_\ell H^1(\bar S,\bar A))^\Gamma\to(M_\ell)_\Gamma\).
Thus taking invariants in the bottom row of 12:77 is not
automatically right exact either.  These are two manifestations of
the same obstruction.

\begin{lemma}[Splitting off divisor classes on a smooth open variety]
\label{lem:B-divisor-splitting}
\status{Proved}
Let \(Y/\F_q\) be smooth, separated and of finite type, and let
\(N_Y\) be the image of \(\BPic(\bar Y)\otimes\Q_\ell\) in
\(H^2(\bar Y,\Q_\ell(1))\), for \(\ell\ne p\).  Then
\(N_Y\) is a \(\Gamma\)-equivariant direct summand.  No
semisimplicity or Tate conjecture for the complementary summand is
assumed.
\end{lemma}
\begin{proof}
First suppose \(Y\) is an open subset of a smooth projective
\(Y'\).  On \(Y'\), the intersection form, formed with an ample
line bundle, is nondegenerate on the rational Neron--Severi group.
Its orthogonal complement therefore splits the divisor-class
subspace in \(H^2(\bar Y',\Q_\ell(1))\).  This is the precise
linear argument in the proof of Yuan's Theorem~1.2; it uses the
Hodge index theorem, not the divisor Tate conjecture~\cite{BYuan2020}.
The kernel of restriction to \(Y\) is spanned by the classes of
the codimension-one components of \(Y'\setminus Y\), by the
Gysin sequence.  It lies inside that divisor summand.  Quotienting
the splitting therefore splits \(N_Y\) inside the image of
restriction; every line bundle on the regular open \(Y\) extends
as a divisor to the regular \(Y'\).  The remaining quotient of
\(H^2(\bar Y,\Q_\ell(1))\) has weights at least one, whereas
the image of restriction has weight zero.  The decomposition into
the relatively prime Frobenius factors of these two weights splits
the image as well.  The Gysin sequence and weight assertion are
given in Balkan--Schreieder's proof of Proposition~4.10 and
Lemma~3.11~\cite{BBalkanSchreieder2026}.

In general an alteration supplies a smooth \(V\), open in a
smooth projective variety, and a proper generically finite
\(\tau:V\to Y\), of degree \(a>0\).  The trace on cohomology
satisfies \(\tau_*\tau^*=a\), and it takes divisor classes to
divisor classes.  These assertions are the projection formula and
cycle compatibility in Balkan--Schreieder's Lemma~2.2.  If \(r_V\)
is the equivariant retraction constructed above, then
\[
 r_Y=a^{-1}\tau_*r_V\tau^*
 :H^2(\bar Y,\Q_\ell(1))\longrightarrow N_Y
\]
restricts to the identity on \(N_Y\).  This proves the result.
The alteration and trace argument works for separated smooth
schemes after taking a proper compactification and a projective
alteration; it does not require \(Y\) itself to be proper.
\end{proof}

\begin{proposition}[The finite-field chain over a regular finite-type base]
\label{prop:B-finite-chain-proper}
\status{Proved}
For the smooth separated base \(S\) above, the natural map
\[
H^1(S,A)[\ell^\infty]\longrightarrow
H^1(\bar S,\bar A)[\ell^\infty]^\Gamma
\]
has finite kernel and cokernel.  Consequently the following four
conditions are equivalent:
\begin{enumerate}
\item \(H^1(S,A)[\ell^\infty]\) is finite;
\item the Kummer map \(A(S)\otimes\Q_\ell\to
H^1_{\Bcts}(S,V_\ell A)\) is an isomorphism;
\item the same map to \(W^\Gamma\) is an isomorphism;
\item \(H^1(\bar S,\bar A)[\ell^\infty]^\Gamma\) is finite.
\end{enumerate}
These are equivalences, without an assumption that their equivalent
finiteness assertions hold.  If \(S\) is projective and \(A\) is
constant, these assertions of finiteness hold unconditionally.
\end{proposition}
\begin{proof}
Put \(B=A^t\), and choose an integer \(a>1\).  Multiplication
\([a]\) on \(B/S\) acts by \(a^j\) on \(R^jf_*\Q_\ell\).
The distinct eigenvalues force all Leray differentials to vanish
and split the degree-two filtration.  The \(a\)-eigenspace of
\(H^2(\bar B,\Q_\ell(1))\) is consequently
\(H^1(\bar S,R^1f_*\Q_\ell(1))=W\).
On relative Picard groups, \([a]^*\) acts by \(a\) on
\(\BPic^0_{B/S}=A\), and by \(a^2\) on relative
Neron--Severi classes; on line bundles pulled back from the base
it acts by one.  The zero section makes
\(\BPic(B)\to\BPic_{B/S}(S)\) surjective.  Thus the
\(a\)-eigenspace of the divisor-class image \(N_B\) is exactly
the Kummer subspace \(M_\ell\) of (B.31).  This identification
also follows at the cocycle level: the line bundle attached to a
section of \(A\) is the pullback of the Poincare bundle, and
choosing its local \(\ell^r\)-division sections gives precisely
the section's Kummer cocycle in \(A[\ell^r]=R^1f_*\mu_{\ell^r}\).

Let \(e_1\) denote the polynomial projector onto the
\(a\)-eigenspace.  Lemma~\ref{lem:B-divisor-splitting} gives a
\(\Gamma\)-equivariant retraction \(r\) onto \(N_B\).
The map \(e_1re_1:W\to M_\ell\) is the identity on
\(M_\ell\), even though \(r\) need not commute with \([a]^*\).
This splits (B.31) equivariantly and makes its invariant sequence
right exact.  All projectors are rational, so the argument also
applies to \(\ell=2\).

Lang's theorem and Lang--Neron identify
\(A(S)\otimes\Q_\ell=(M_\ell)^\Gamma\).
The kernel in continuous Hochschild--Serre is
\(H^0(\bar S,V_\ell A)_\Gamma=0\).  Indeed evaluation at a
closed point of \(S\) embeds this global-section space in the
Tate module of that abelian fibre; a power of Frobenius on it has
weight \(-1\), excluding eigenvalue one.  The diagram of global
Kummer and the now exact invariant sequence of (B.31) therefore
identifies
\(V_\ell H^1(S,A)\) with
\((V_\ell H^1(\bar S,\bar A))^\Gamma\).
The finite sheaf \(A[\ell]\) and Kummer prove cofiniteness of
both torsion groups.  Their natural comparison consequently has
finite kernel and cokernel, and all four conditions are equivalent.
This proves, in particular, that the missing transgression in
(B.29) has finite image in the stated regular-base setting.

For constant \(A=A_0\times S\), sections modulo the finite
constant group identify, up to a finite ambiguity when choosing an
Albanese base point, with \(\Hom(\operatorname{Alb}_S,A_0)\).
Moreover
\(W=\Hom(V_\ell\operatorname{Alb}_S,V_\ell A_0)\).
Tate's homomorphism theorem for finite-field abelian varieties
identifies its invariants with those homomorphisms tensored with
\(\Q_\ell\).  Condition (3) follows~\cite[Main Theorem, \S1]{BTate1966}.
\end{proof}

Keller's Theorem~2.7.19 and Remark~2.7.20(d) identify the same
proper-base finiteness assertion with equality of algebraic and
analytic rank.  The height pairing is nondegenerate in the required
case; it cannot be silently omitted from the theorem's printed
condition before that remark is applied~\cite{BKeller2019}.
Thus 12:80 is a proved equivalence with BSD in this scope, not a
proof of BSD for an arbitrary nonconstant abelian scheme.

\begin{proposition}[Relative curves without a section]
\label{prop:B-relative-curve}
\status{Proved}
Let \(S/\F_q\) be smooth, separated and connected of finite type,
and let \(f:Y\to S\) be a smooth proper curve with geometrically
connected fibres and relative Jacobian \(J\).  For \(\ell\ne p\),
\[
 V_\ell\bigl(\Br(Y)/\operatorname{im}\Br(S)\bigr)
 \simeq V_\ell H^1(S,J).
 \tag{B.32}
\]
The kernel of \(\Br(S)[\ell^\infty]\to\Br(Y)[\ell^\infty]\)
is finite.  Consequently the following are equivalent:
\begin{enumerate}
\item \(\Br(S)\to\Br(Y)\) is surjective modulo a finite
\(\ell\)-primary group;
\item \(H^1(S,J)[\ell^\infty]\) is finite;
\item \(H^1(\bar S,\bar J)[\ell^\infty]^\Gamma\) is finite.
\end{enumerate}
If \(f\) has a section, the natural edge map gives the stronger
integral isomorphism
\[
 (\Br(Y)/\Br(S))[\ell^\infty]
 \simeq H^1(S,J)[\ell^\infty].
\]
\end{proposition}
\begin{proof}
Choose a closed point of degree \(d>0\) on the generic fibre.
Its closure in the regular scheme \(Y\) is a Cartier divisor;
its line bundle \(L\) has degree \(d\) on every geometric
fibre, because degree in a proper smooth family of curves is locally
constant.  The closure is not asserted to be finite over \(S\).
Cup product with \(c_1(L)\) induces the isomorphism
\(R^0f_*\Q_\ell\to R^2f_*\Q_\ell(1)\), namely
multiplication by \(d\).  Deligne's Proposition~2.4, equation
(2.5), therefore applies directly; for a curve this is the only
nontrivial Lefschetz condition.  In particular its Leray filtration
in degree two is split over \(\Q_\ell\), with successive terms
\[
 H^2(S,\Q_\ell(1)),\qquad
 H^1(S,V_\ell J),\qquad \Q_\ell.
\]
The first term is the image of \(f^*\), and the last term is
split by \(c_1(L)/d\)
\cite[Proposition~2.4, (2.5), pp.~111--112]{c_deligne1968}.

Put \(P=\BPic_{Y/S}=R^1f_*\Gm\).  Etale locally the
curve has a section, so there is an exact sequence of sheaves
\(0\to J\to P\xrightarrow{\deg}\Z\to0\).
The image of \(L\) has degree \(d\), hence the boundary
\(\Z\to H^1(S,J)\) has finite cyclic image of order dividing
\(d\).  Moreover \(H^1(S,\Z)=0\).  Thus
\(H^1(S,J)\to H^1(S,P)\) has finite kernel and is surjective.
The obstruction to lifting a section of \(P\) to a line bundle
on \(Y\) is in
\(\ker(\Br(S)\to\Br(Y))\).  This kernel is killed by
\(d\): restrict to the function field of the regular integral
\(S\), then to the chosen degree-\(d\) point, and apply
Brauer corestriction.  Injectivity of \(\Br(S)\to\Br(k(S))\)
returns the assertion to \(S\).  Its \(\ell\)-primary part
is finite, because Kummer and finiteness of
\(H^2(S,\mu_\ell)\) make \(\Br(S)[\ell^\infty]\)
cofinitely generated.

It remains to identify the middle algebraic subspace, rather than
infer anything from the rank of a Leray term alone.  The relative
Chern class of a degree-zero relative line bundle is its Kummer
class in \(H^1(S,V_\ell J)\): this follows by choosing etale
local lifts and \(\ell^r\)-division sections and comparing the
resulting transition cocycles.  The bounded Brauer obstruction just
proved permits every section of \(J\) to be lifted after
multiplication by \(d\).  The divisor-class image in the
middle term is therefore exactly \(J(S)\otimes\Q_\ell\).
In the first term it is the divisor-class image from \(S\),
and in the last term it is all of \(\Q_\ell\), generated by
\(L\).  Equivalently, one can apply this argument to the
\(\ell\)-adic completions of the Picard groups and then tensor
with \(\Q_\ell\); finite kernels and the bounded obstruction
contribute no term.  Quotienting the three Leray terms by these
Chern-class images, the two Kummer sequences give
\[
 \frac{V_\ell\Br(Y)}{V_\ell\Br(S)}
 \simeq
 \frac{H^1(S,V_\ell J)}{J(S)\otimes\Q_\ell}
 \simeq V_\ell H^1(S,J).
\]
This proves (B.32).  All torsion groups in this comparison are
cofinitely generated on their \(\ell\)-primary parts, again by
Kummer and finite etale cohomology.  Thus vanishing of these rational
Tate modules is equivalent to finiteness, proving (1) equivalent
to (2); Proposition~\ref{prop:B-finite-chain-proper} gives (3).

When there is a section, \(P=J\times\Z\).  On strict local
bases the curve has a section, and proper base change and Kummer
show \((R^2f_*\Gm)[\ell^\infty]=0\).  The resulting
Leray edge sequence on primary torsion has target
\(H^1(S,P)\).  Its next differential lands in the kernel of
\(H^3(S,\Gm)\to H^3(Y,\Gm)\), which is zero because
restriction to the section splits pullback.  This proves the final
integral assertion without assuming \(H^3(S,\Gm)=0\).
\end{proof}

\begin{remark}[A source error avoided in the preceding proof]\label{rem:B-Keller-closure-error}
\status{Refuted; corrected statement proved}
Keller's Lemma~2.17 claims that the closure of a generic-fibre point
of a proper relative curve is finite over an integral base.  This
is false in the printed generality.  Over any finite field take
\(S=\mathbb A^2_{x,y}\), \(Y=\mathbb P^1_{[u:v]}\times S\),
and the generic point \([u:v]=[x:y]\).  Its closure is
\(uy-vx=0\), the blowup of the origin, whose fibre over
\((0,0)\) is the whole \(\mathbb P^1\).  On its two standard
charts the coordinate rings are \(k[x,t]\), \(y=xt\), and
\(k[y,s]\), \(x=ys\); it is integral and has no additional
irreducible component.  This explicitly contradicts the lemma's
quasi-finiteness step.  The degree of the associated line bundle,
which is all the proof above uses, remains one on every fibre.
Thus the section-free part of the preceding proposition is not
inherited from that lemma or from its use in Corollary~2.18 and
Theorem~4.27~\cite[Lemma~2.17]{BKeller2016}.
\end{remark}

The proposition proves the ``via des jacobiennes'' equivalence on
12:79 over its stated regular base, without adding a section
hypothesis.  It does not prove the equivalent finiteness assertions
for every nonconstant Jacobian.  When \(S\) is a smooth projective
curve, \(\Br(S)=0\), and the assertion is precisely finiteness
of the Brauer group of the surface \(Y\), equivalently its divisor
Tate conjecture by Yuan's Theorems~1.2 and 1.3(4).

For a nonproper smooth curve base, Qin's Propositions~2.4--2.5 and
Lemma~2.6 provide an independent proof of the comparison, including
possible bad reduction at the omitted points~\cite{BQin2024}.
The argument above controls higher-dimensional nonproper regular
bases too, by splitting only the algebraic subspace, a weaker and
proved property than semisimplicity of all of \(W\).
\paragraph{The chain over every finite-type base.}
The smooth-base splitting above is enough to remove regularity from
the Kummer comparison itself.  One needs injectivity on sections,
not descent of a family of cohomological retractions along a
hypercover.

\begin{lemma}[Sections and their completion on a reduced base]
\label{lem:B-sections-reduced}
\status{Proved}
Let \(S/\F_q\) be reduced and of finite type, without a
separatedness assumption, and let \(A/S\) be an abelian scheme.
Put \(C_S=A(\bar S)\).  Then \(C_S/(C_S)_{\rm tors}\)
is finitely generated, its Galois action has finite image, and
the natural map gives an isomorphism
\[
 C_S\otimes_{\Z}\Q_\ell
 \xrightarrow{\sim}
 (C_S)^{\wedge}_\ell\otimes_{\Z_\ell}\Q_\ell
 \qquad (\ell\ne p).
\]
There exists a finite disjoint union \(V\) of smooth separated
finite-type \(\F_q\)-schemes with a jointly surjective map
\(v:V\to S\).  Pullback embeds \(C_S\) in \(C_V\)
and embeds their quotients by torsion.
\end{lemma}
\begin{proof}
Take a finite affine cover of \(S\).  Alter each reduced
irreducible component of each member of the cover.  De Jong's
Theorem~4.1 and Remark~4.2 give a smooth separated alteration,
since \(\F_q\) is perfect; their disjoint union is the
required \(V\)~\cite{c_dejong1996}.  The construction and
surjectivity persist after extension to \(\overline{\F}_q\).
If a section becomes zero on \(\bar V\), its equalizer with
the zero section is a closed subscheme of \(\bar S\) whose
underlying set is all of \(\bar S\).  Its ideal is therefore
contained in the nilradical and is zero.  Thus pullback is
injective.  Applying the same argument to an integer multiple of
a section shows injectivity modulo torsion as well.

For each connected component of \(\bar V\), restriction to
its generic point embeds its sections into the points of an
abelian variety over a finitely generated regular extension of
\(\overline{\F}_q\).  Lang--Neron, in the exact form of
Conrad's Theorems~2.1 and~7.1, says that the quotient by the
constant trace is finitely generated~\cite{BConrad2006}.  All
\(\overline{\F}_q\)-points of an abelian variety over
\(\overline{\F}_q\) are torsion: the variety and each point
descend to a finite field.  The quotient of its function-field
points by all torsion is consequently finitely generated.  A
subgroup of a finitely generated abelian group is finitely
generated.  Applying this twice, first to \(C_V\) and then to
\(C_S\), proves the assertion about sections modulo torsion.
Representatives of finitely many generators are defined over one
finite extension of \(\F_q\), which proves the finite-image
assertion.

Evaluation at one geometric point of each connected component of
\(\bar S\) injects \(C_S[\ell^\infty]\) into a finite
sum of groups \(A_{\bar s}[\ell^\infty]\).  Indeed a section
of a finite etale group scheme which vanishes at one point of a
connected scheme vanishes everywhere.  Thus the Pontryagin dual
of \(C_S[\ell^\infty]\) is a quotient of a finite free
\(\Z_\ell\)-module; it is finitely generated.  Equivalently,
this torsion group is a sum of finitely many copies of
\(\Q_\ell/\Z_\ell\) and a finite group.  Its
\(\ell\)-adic completion has zero rationalization.  The
prime-to-\(\ell\) torsion is \(\ell\)-divisible and has
zero completion.  Finally the extension of the torsion subgroup
by the finite free group \(C_S/(C_S)_{\rm tors}\) splits as
an extension of abelian groups.  Computing completions in that
splitting proves the displayed natural isomorphism; its statement
does not depend on the splitting chosen.
\end{proof}

\begin{theorem}[The finite-field chain for arbitrary finite-type schemes]
\label{thm:B-finite-chain-all-bases}
\status{Proved}
Let \(S/\F_q\) be any scheme of finite type, let \(A/S\)
be an abelian scheme, and let \(\ell\ne p\), including
\(\ell=2\) when \(p\ne2\).  No regularity, reducedness,
properness or separatedness of \(S\) is required.  Set
\[
 W_S=H^1_{\Bcts}(\bar S,V_\ell A),\qquad
 N_S=A(\bar S)\otimes\Q_\ell.
\]
The geometric Kummer sequence
\[
 0\longrightarrow N_S\longrightarrow W_S
 \longrightarrow V_\ell H^1_{\et}(\bar S,\bar A)
 \longrightarrow0
\]
has a \(\Gamma\)-equivariant splitting.  The natural map
\[
 H^1_{\et}(S,A)[\ell^\infty]\longrightarrow
 H^1_{\et}(\bar S,\bar A)[\ell^\infty]^\Gamma
\]
has finite kernel and cokernel.  The four conditions in
Proposition~\ref{prop:B-finite-chain-proper} are therefore
equivalent for every such \(S\).
\end{theorem}
\begin{proof}
First suppose that \(S\) is reduced and take \(V\) as in
Lemma~\ref{lem:B-sections-reduced}.  Kummer maps commute with
pullback.  Write \(i:N_S\hookrightarrow N_V\) for the
injection of that lemma and \(v^*:W_S\to W_V\) for
cohomological pullback.  The proof of
Proposition~\ref{prop:B-finite-chain-proper}, component by
component, supplies a \(\Gamma\)-equivariant retraction
\(r_V:W_V\to N_V\).  The action of \(\Gamma\) on
\(N_V\) factors through a finite group \(G\).  Choose any
linear retraction \(b:N_V\to N_S\) of \(i\) and average:
\[
 b_G=\frac1{|G|}\sum_{g\in G}g\,b\,g^{-1}.
\]
The scalar \(|G|^{-1}\) is available in \(\Q_\ell\),
even when \(\ell\mid |G|\).  Hence
\[
 r_S=b_G\,r_V\,v^*:W_S\longrightarrow N_S
 \quad\hbox{satisfies}\quad r_S|_{N_S}=1.
\]
This proves the splitting.  It does not require \(v^*\) to be
injective on the whole of \(W_S\), nor compatible retractions
on intersections of the covering alterations.

For a nonreduced \(S\), filter its nilpotent ideal by powers,
so that successive closed immersions have square-zero ideals
\(I\).  On the small etale site, smoothness of \(A\) gives
a surjection from the sheaf of its sections to the sheaf on the
thickening's reduction.  Its kernel for a square-zero step is
\(\operatorname{Lie}(A)\otimes I\), an additive sheaf killed
by \(p\).  The cohomology exact sequence and induction show
that the kernels and cokernels on \(H^0\) and \(H^1\) are
of bounded \(p\)-power exponent.  Such maps are isomorphisms
on \(\ell\)-primary torsion and after tensoring sections
with \(\Q_\ell\).  The finite etale sheaves
\(A[\ell^a]\) and their cohomology are unchanged by a
nilpotent thickening.  Thus the Kummer sequence and its splitting
reduce to those already proved for \(S_{\rm red}\).

It remains to verify the passage from geometric splitting to the
stated finite comparison.  Norms from finite constant extensions
give
\[
 A(S)\otimes\Q_\ell=N_S^\Gamma.
\]
Indeed a finite set of section representatives is defined over
one finite extension, and averaging over that extension descends
the invariant rational combination; injectivity is immediate.
Furthermore
\(H^0(\bar S,V_\ell A)_\Gamma=0\).  To see this,
evaluate at a finite set of geometric points, one on each
connected component, lying over closed points of \(S\).
This is injective on sections of a lisse sheaf.  After a finite
constant extension it is equivariant into a finite sum of Tate
modules of abelian varieties over finite fields.  Their
geometric Frobenius eigenvalues have weight \(-1\), so no
positive power has eigenvalue one.  Continuous
Hochschild--Serre therefore identifies
\[
 H^1_{\Bcts}(S,V_\ell A)=W_S^\Gamma.
\]
Taking invariants in the split geometric Kummer sequence and
comparing with global Kummer gives the natural isomorphism
\[
 V_\ell H^1_{\et}(S,A)
 \xrightarrow{\sim}
 \bigl(V_\ell H^1_{\et}(\bar S,\bar A)\bigr)^\Gamma.
\]
The group \(H^1(S,A)[\ell]\) is a quotient of the finite
group \(H^1(S,A[\ell])\), and the same holds over
\(\overline{\F}_q\).  Finite-coefficient etale finiteness
also holds for a nonseparated finite-type scheme by a finite
affine cover and its Cech spectral sequence.  Consequently both
\(\ell\)-primary groups are cofinitely generated.  Taking
invariants preserves this property, and the Tate module of the
invariants is the invariant Tate module, since invariants commute
with inverse limits.  A map between cofinitely generated
\(\ell\)-primary torsion groups with an isomorphism on
rational Tate modules has finite kernel and cokernel.  This proves
the comparison and all four equivalences.
\end{proof}

The transcription of 12:79 itself explicitly retains the condition
that the base be regular.  The relative-curve proposition above
therefore proves the requested Jacobian--Brauer equivalence in
that stated scope.  The present theorem removes the former
singular-base gap in the Kummer chain; it does not prove its
equivalent finiteness conditions for every nonconstant abelian
scheme.


\begin{example}[Why the positive-degree argument retains its regularity hypothesis]
\label{ex:B-singular-infinite-period}
\status{Proved counterexample to an extension beyond the stated manuscript scope}
There is a reduced projective surface \(S/\F_2\) and a smooth
proper genus-one scheme \(f:X\to S\) for which every global line
bundle has relative degree zero.  Thus the positive-degree class
used in Proposition~\ref{prop:B-relative-curve} cannot be asserted
on an arbitrary singular base.  This example does not assert that
the Brauer quotient is infinite.

Take
\[
 E:\ y^2z+yz^2+x^3+x^2z=0\ \subset T=\mathbf P^2_{\F_2},
 \qquad O=[0:1:0],
 \qquad S=T\amalg_E T.
\]
The derivatives of the cubic are \(x^2,z^2,y^2+x^2\), with no
common projective zero, so \(E\) is an elliptic curve.  The
pushout is the explicit finite \(T\)-scheme
\[
 S=\operatorname{Spec}_T(\mathcal O_T\oplus\mathcal O_T(-3)),
 \qquad
 R[t]/(t(t-f_E))\simeq R\times_{R/(f_E)}R
\]
on each trivializing affine, where multiplication on the second
summand uses the cubic section \(f_E\).  It is reduced because
the displayed algebra embeds in the product of two domains, and
projective because it is finite over \(T\).

Let \(A=E\times S\).  The closed-cover sequence on the small
etale site is
\[
 0\longrightarrow A\longrightarrow i_{1*}A_T\oplus i_{2*}A_T
 \xrightarrow{\mathrm{res}_1-\mathrm{res}_2}i_{E*}A_E
 \longrightarrow0.
\]
For surjectivity it suffices to pass to a strict henselian local
ring.  A section into the smooth curve \(E\) on its closed
quotient lifts across the resulting henselian pair; see the Stacks
Project, Lemma~15.13.3~\cite{BStacksSmoothLift}.  The kernel is
the fibre-product description of \(S\).  Every map
\(\mathbf P^2\to E\) is constant, since its restriction to
each projective line is constant, whereas the rigidity lemma gives
\[
 A(T)=E(\F_2),\qquad
 A(E)=E(\F_2)\oplus\operatorname{End}_{\F_2}(E).
\]
The cohomology boundary therefore injects
\(\operatorname{End}_{\F_2}(E)\) into \(H^1_{\et}(S,A)\).
In particular \(\alpha=\delta(1_E)\) has infinite order:
\([m]:E\to E\) has degree \(m^2\) and is nonconstant for every
nonzero integer \(m\).

To represent this class by a scheme, glue two copies
\(X_i=E\times T\) along \(E\times E\), using the identity
on one side and \((a,t)\mapsto(a+t,t)\) on the other.
Both attaching maps are closed immersions, so the pushout is a
scheme by the Stacks Project, Lemma~37.67.8; Lemma~37.67.9(4)
gives smoothness over \(S\)~\cite{BStacksClosedPushout,BStacksPushoutMorphisms}.  The same local lifting just used
trivializes the translation cocycle on an etale cover.  Thus this
scheme is the \(A\)-torsor of class \(\pm\alpha\), locally
isomorphic to \(E\times S\), and is proper with geometrically
connected genus-one fibres.

Under the principal polarization \(E\simeq E^\vee\), translation
by \(a\) acts on \(\operatorname{Pic}^d(E)\) by translation
by \(\pm da\): represent a divisor by its sum in \(E\).
Consequently \(\operatorname{Pic}^d_{X/S}\) is a torsor of
class \(\pm d\alpha\).  A global line bundle of degree \(d\)
would give a section of that torsor, forcing \(d\alpha=0\),
hence \(d=0\).  This proves the assertion, and also shows that
\(X\to S\) is not projective.  The cubic smoothness, its five
\(\F_2\)-points, and the displayed fibre-product algebra were
checked exactly by
\texttt{part\_b/code/singular\_genus\_one\_checks.py}; its output
is preserved with the computations.  The infinite-order and torsor
claims are proved above, not inferred from the finite computation.
\end{example}

\begin{example}[A fully specified nonzero degree-two group on 12:78]\label{ex:B-H2-constant}
\status{Proved}
Let \(C/\F_2\) be the good reduction of the displayed
\(X_0(11)\) model, and let \(A=C\times C\to C\) be its
constant Jacobian scheme.  Take \(\ell=3\).  Direct enumeration
gives \(\#C(\F_2)=5\), so Frobenius on \(V_3C\) has
polynomial \(t^2+2t+2\).  Its two roots are distinct, and the
centralizer of a companion matrix has dimension two.  Kummer,
the constant-case result just proved, and Hochschild--Serre give
\[
 V_3H^2_{\et}(C,A)
 \simeq H^2_{\Bcts}(C,V_3A)
 \simeq H^1_{\Bcts}(\bar C,V_3A)_\Gamma
 \simeq\End_{\F_2}(C)\otimes\Q_3,
\]
of dimension two.  The geometric degree-two term has weight one,
so contributes no invariants; in the degree-one term,
\(H^1(\bar C,\Q_3)\otimes V_3C=(V_3C)^\vee\otimes V_3C\).
Semisimplicity identifies its coinvariants with its invariants,
and Tate's Main Theorem identifies the latter with endomorphisms.
The point count and exact centralizer calculation are reproduced in
\texttt{part\_b/code/finite\_field\_checks.py}.

This verifies the intended nonvanishing in R6.  If a dual Tate module
is replaced in the transcription by the Tate module of the dual
abelian variety, a twist must be retained:
\((V_\ell J)^\vee=V_\ell(J^t)(-1)\).
An integral isomorphism also cannot be read off from a displayed
\(\Q_\ell\)-vector-space calculation without tracking finite
Kummer kernels.
\end{example}

\subsection{B.4: the rational number-ring calculation and the integral correction}
\begin{theorem}[All rational twists on the compactified integer ring]\label{thm:B-SpecZ-Q}
\status{Proved}
In the Flach--Morin theory of \(\overline{\Spec\Z}\),
\[
H^i_{\BW}(\overline{\Spec\Z},\Q(n))=
\begin{cases}
\Q,&(n,i)=(0,0)\text{ or }(1,3),\\
\Q,&n\ge3\text{ odd},\ i=1,\\
\Q,&n\le-2\text{ even},\ i=2,\\
0,&\text{otherwise}.
\end{cases}\tag{B.5}
\]
The unit in degree zero pairs perfectly with the fundamental class in
degree three and twist one.  For every odd \(n\ge3\), the degree-one
regulator line in twist \(n\) pairs perfectly with the degree-two line
in twist \(1-n\).  The dualizing object is \(\Q(1)[3]\).
\end{theorem}
\begin{proof}
Flach--Morin's Corollary~3.12, in absolute dimension one, is
\[
\BRGamma_{\BW}(\bar X,\Q(n))\simeq
\BRGamma_{\BM}(X,\Q(n))\oplus
\BRHom(\BRGamma_{\BM}(X,\Q(1-n)),\Q[-3]).\tag{B.6}
\]
For \(X=\Spec\Z\), ordinary rational motivic cohomology is \(\Q\)
in degree zero and twist zero.  In positive twists its only possible
nonzero groups are \(H^1_{\BM}(X,\Q(n))\): twist one is
\(\Z^\times\otimes\Q=0\), while Borel's ranks give a one-dimensional
space for odd \(n\ge3\), and zero for even \(n\ge2\).
All negative-twist groups vanish.  Indeed, rational motivic cohomology
of this regular scheme is its Adams-graded rational \(K\)-theory;
the finite-field localization terms are torsion, the positive
even \(K\)-groups have rank zero, and the odd ranks are Borel's
Proposition~12.2~\cite{BBorel1974}.  Substitution in (B.6) proves
every case of (B.5).  Theorem~3.22 of Flach--Morin identifies the
pairing on (B.6) with the two cross-evaluation maps, so it is perfect
and has the claimed degree and twist~\cite[Corollary~3.12,
Theorem~3.22, and \S5.8.3]{BFlachMorin2018}.
\end{proof}

Theorem~\ref{thm:B-primitive-generator} below gives a terminating
construction of a primitive integral class in each odd positive line.
Its evaluation-dual gives the negative rational line.  These choices
are specified by finite searches; they are not canonical motivic
normalizations.

\begin{proposition}[The omitted order-two groups]\label{prop:B-integral-errors}
\status{Refuted; corrected statement proved}
The two integral tables asserted in 15:173 are false in Flach--Morin
Weil--etale cohomology.  Their correct values are
\[
H^i_{\BW}(\overline{\Spec\Z},\Z(1))=
\begin{cases}\Z/2,&i=1,\\\Z,&i=3,\\0,&\text{otherwise},\end{cases}
\qquad
H^i_{\BW}(\overline{\Spec\Z},\Z)=
\begin{cases}\Z,&i=0,\\\Z/2,&i=3,\\0,&\text{otherwise}.\end{cases}\tag{B.7}
\]
\end{proposition}
\begin{proof}
The number-ring table in Flach--Morin \S5.8.3 gives, in twist one,
the units, the class group and \(\Z\) in degrees one, two and three.
For \(\Z\) these are \(\{\pm1\}\), zero and \(\Z\).
In twist zero it gives \(\Z\) in degree zero,
\(\Hom(\mathcal O_K^\times,\Z)\oplus\BCl(\mathcal O_K)^D\)
in degree two, and \(((\mathcal O_K^\times)_{\rm tors})^D\) in degree
three, where \(D=\Hom(-,\Q/\Z)\).  Specializing to \(K=\Q\)
proves (B.7).  The order-two classes are dual under the integral derived
pairing; tensoring with \(\Q\) removes both.  This checks, and repairs,
the exact R5 step used in Theorem~\ref{thm:B-SpecZ-Q}.
\end{proof}

\paragraph{Effective completion of the integral calculation.}
The source's structural table alone leaves finite arithmetic groups
and primitive free classes unspecified.  The continuation below fills
both requirements: Theorem~\ref{thm:B-effective-odd-torsion} computes
all odd-primary invariant factors with cohomological generators,
Proposition~\ref{prop:B-two-primary} supplies every two-primary
factor, and Theorem~\ref{thm:B-primitive-generator} constructs
primitive integral free classes by terminating finite searches.
This is an effective presentation for each requested twist, with no
claim of an efficient bound or of uniform Vandiver vanishing.

\begin{proposition}[The complete structural table]
\label{prop:B-integral-structure}
\status{Proved}
For \(m\ge2\), put
\(C_m=H^2_{\et}(\Spec\Z,\Z(m))\), a finite group, and
\(a_m=1\) for odd \(m\ge3\), \(a_m=0\) for even \(m\).
Let \(w_m\) be the order of
\(H^0(\Q,\Q/\Z(m))\).  Explicitly,
\[
 w_m=\begin{cases}
 2,&m\text{ odd},\\
 2^{2+v_2(m)}\displaystyle\prod_{\substack{p\text{ odd prime}\\p-1\mid m}}
 p^{1+v_p(m)},&m\text{ even}.
 \end{cases}
\]
Define \(\epsilon_{i,n}\) to be one if
\(i\equiv n\pmod2\) and either \(1\le i\le n\) or
\(n<i<0\), and zero otherwise.  In addition to (B.7), all the
integral groups have the following structure:
\[
H^i_{\BW}(\overline{\Spec\Z},\Z(n))\simeq
\begin{cases}
\Z^{a_n}\oplus\Z/w_n,&n>1,\ i=1,\\
C_n,&n>1,\ i=2,\\
(\Z/2)^{\epsilon_{i,n}},&n>1,\ i\ge3,\\
(\Z/2)^{\epsilon_{i-1,n}},&n<0,\ i<0,\\
\Z^{a_{1-n}}\oplus C_{1-n}^{D},&n<0,\ i=2,\\
\Z/w_{1-n},&n<0,\ i=3,\\
0,&\text{all remaining cases with }n\notin\{0,1\}.
\end{cases}\tag{B.22}
\]
The direct-sum splittings and generators of the free groups in this
display are not canonical.  The arithmetic computation denoted by \(C_m\) is explicitly
supplied by the two following results.
\end{proposition}
\begin{proof}
The table in Flach--Morin \S5.8.3 expresses the answer through
\[
 H^{1,m}=H^1_{\et}(\Spec\Z,\Z(m)),
\]
the groups \(C_m\), their
integral and Pontryagin duals, and the displayed real-place terms.
Its finite generation and rank calculation give
\(H^{1,m}\simeq\Z^{a_m}\oplus(H^{1,m})_{\rm tors}\).
The rational-to-torsion coefficient sequence, together with the
positive-twist localization sequence, identifies the torsion subgroup
with \(H^0(\Q,\Q/\Z(m))\).  For each odd prime \(p\), the
cyclotomic character has image \(\Z_p^\times\).  Its \(m\)-th
power acts trivially on a cyclic subgroup of order \(p^b\) exactly
when \(p-1\mid m\) and \(b\le1+v_p(m)\).  At two, use
\(\Z_2^\times=\{\pm1\}\times(1+4\Z_2)\): an odd twist
has invariants of order two, and an even twist has invariants of
order \(2^{2+v_2(m)}\).  This proves the formula for \(w_m\).
Substitution in the source's integral table proves (B.22), including
the bounded negative-degree terms.  Rationalization recovers
Theorem~\ref{thm:B-SpecZ-Q}.
\end{proof}

\begin{proposition}[Every two-primary entry and its generator]
\label{prop:B-two-primary}
\status{Proved}
For the finite group \(C_m\) in (B.22),
\[
 C_m[2^\infty]=
 \begin{cases}\Z/2,&m\ge2\text{ even},\\0,&m\ge3\text{ odd}.
 \end{cases}
\]
For even \(m\), a generator is the inverse image, under the
canonical localization and completion isomorphism, of
\[
 \delta\bigl(\iota_m([-1])\bigr)
 \in H^2(\Spec\Z[1/2],\Z_2(m)).
\]
Here \([-1]\in H^1(\Spec\Z[1/2],\Z/2)\) is the
Kummer class of \(-1\), \(\iota_m:\Z/2\to\Q_2/\Z_2(m)\)
is the order-two inclusion, and \(\delta\) is the connecting map
for \(0\to\Z_2(m)\to\Q_2(m)\to\Q_2/\Z_2(m)\to0\).
\end{proposition}
\begin{proof}
Rognes--Weibel's Lemma~5.5 identifies
\(H^1(\Z[1/2],\Q_2/\Z_2(m))\) with its real restriction
\(\Z/2\) for nonzero even \(m\); its proof gives precisely
\(\iota_m([-1])\) as generator.  For positive even \(m\),
Borel--Soule gives \(H^1(\Z[1/2],\Q_2(m))=0\) and
\(H^2(\Z[1/2],\Q_2(m))=0\).  Hence \(\delta\) is an
isomorphism.  For odd \(m\ge3\), Proposition~6.13 of the same
source says that the finite group
\(H^2(\Z[1/2],\Z_2(m))\) has two-rank \(s+t-1\),
where \(s\) counts primes above two and \(t\) is the
class group's two-rank.  For \(\Q\), \(s=1,t=0\), so this
group vanishes~\cite[Lemma~5.5, Propositions~6.12--6.13]{BRognesWeibel2000}.
Finally, the localization contribution at two is in the positive
finite-field twist \(m-1\), of odd order \(2^{m-1}-1\).
Thus localization and etale motivic comparison identify the just
computed groups with \(C_m[2^\infty]\).  In particular these
are the unmodified motivic groups; the separate real-place terms in
(B.22) must still be retained.
\end{proof}

\subsubsection{An effective finite presentation for the odd-primary terms}
\begin{theorem}[Finite arithmetic computation of each odd-primary term]
\label{thm:B-effective-odd-torsion}
\status{Proved as a terminating effective presentation}
\status{Proved}
Fix an integer \(n\ge2\).  There is an unconditional terminating
algorithm, specified below using integer polynomial enumeration and
finite matrices, which returns the invariant factors and cohomological
generators of the odd-primary subgroup of
\(C_n=H^2_{\et}(\Spec\Z,\Z(n))\).  It requires neither Vandiver's
conjecture nor a class-group vanishing assertion.  This is an effective
presentation, not a claim that the resulting invariant factors obey
one presently proved closed vanishing pattern.
\end{theorem}
\begin{proof}
We give the complete finite procedure and its termination argument.
Write \(m=2n-2\), \(N=2m+1\), and put
\[
 A_N=2^{4N}(N+1)^{2N^2},\qquad
 w_N=N(N+1)/2-N,
\]
\[
 K_n=A_N^{(N+1)\binom{A_N}{N+w_N-m-1}},\qquad
 B_n=\max(4n-2,K_n).
\tag{B.A1}
\]
These are explicit positive integers.  Only odd primes \(\ell\le B_n\)
need be considered.  Indeed, Soul\'e's Theorem~3 gives
\[
 \prod_{\ell\ge4n-1}
 \#H^2(\Spec\Z[1/\ell],\Z_\ell(n))\le k(2n-2).
\]
In his Proposition~3 and Theorems~1--2, the bound for \(\Q\) is
\(k(m)=a(N)^{(N+1)\binom{a(N)}{N+w_N-m-1}}\).
His Proposition~2 and the calculation in \S4.4 give
\(a(N)=2^{4N}(\prod_iT_i)^{1/2}\), with
\(T_i=\gamma^N\prod_{j\ne i}B_j^2\),
\(B_j=(j+1)/2\), and \(\gamma\le1+N/4\).
Thus \(a(N)\le A_N\), and the binomial estimate is monotone in this
range, giving \(k(m)\le K_n\).  A nonzero \(\ell\)-group has order
at least \(\ell\); hence every prime greater than \(B_n\) contributes
zero.  We use only the explicit bound, not an asymptotic or a
conjectural size estimate~\cite[Proposition~3, Theorems~1--3,
and \S4.4]{Soule2003Bound}.

Fix one of these primes, an integer \(r\ge1\), and set
\(q=\ell^r\), \(D=q\varphi(q)\), \(U=\Spec\Z[1/\ell]\).
Let \(\mathcal G=\operatorname{Gal}(\Q_{\{\ell,\infty\}}/\Q)\)
be the Galois group of the maximal extension unramified at every
finite prime other than \(\ell\).  On the cyclic module
\(M_q=\Z/q(1-n)\), let
\(a_g=\chi_q(g)^{1-n}\in(\Z/q)^\times\).
An inhomogeneous continuous crossed cocycle
\(z:\mathcal G\to M_q\) defines the homomorphism
\[
 g\longmapsto(z(g),\chi_q(g)):
 \mathcal G\longrightarrow
 (\Z/q)\rtimes(\Z/q)^\times,
\tag{B.A2}
\]
where the action on the first factor is the \((1-n)\)-th power.
Its image has order at most \(D\).  Consequently every such cocycle
factors through a Galois extension of \(\Q\) of degree at most
\(D\), ramified at no finite prime except \(\ell\).

Here is a finite, explicitly bounded enumeration of all those fields.
Put
\[
 \Delta=\ell^{D(1+\lfloor\log_\ell D\rfloor)}.
\tag{B.A3}
\]
A Galois field \(E/\Q\) of degree \(d\le D\), ramified only at
\(\ell\), has \(|\operatorname{disc}E|\le\Delta\).
To verify this, at \(\ell\) write the ramification and residue
degrees as \(e,f\).  Over the maximal unramified subextension an
Eisenstein polynomial for a uniformizer has degree \(e\).  The
valuations of its nonzero derivative terms have distinct residues
modulo \(e\).  Its different exponent is therefore at most
\(e-1+e v_\ell(e)\), the valuation of the leading derivative term.
Taking the norm of the different at all primes gives
\[
 v_\ell(\operatorname{disc}E)
 \le d(1+v_\ell(e))\le D(1+\lfloor\log_\ell D\rfloor).
\]
There are no other finite discriminant factors.

For each \(2\le d\le D\), enumerate all monic integer
polynomials of degree \(d\) whose coefficients have absolute value
at most
\[
 H(d,\Delta)=2^d\bigl(d^{d+3}\Delta\bigr)^d.
\tag{B.A4}
\]
Every degree-\(d\) field with discriminant bounded by \(\Delta\)
occurs.  In fact, in its standard real Minkowski embedding the
integer lattice has covolume at most \(\sqrt\Delta\), and every
nonzero vector has Euclidean norm at least one: an algebraic integer
of nonzero norm has at least one conjugate of absolute value at
least one.  The unit ball contains the cube of side \(2/\sqrt d\).
Minkowski's second theorem therefore bounds the last successive
minimum by \(d^{d/2}\sqrt\Delta\).  Choose \(d\) linearly
independent integers \(b_i\) of norm at most
\(2d^d\sqrt\Delta\).  For the \(t=d(d-1)/2\) pairs of
distinct embeddings, each equality
\(\sigma(\sum c_i b_i)=\tau(\sum c_i b_i)\) excludes at most
\((t+1)^{d-1}\) elements of the grid
\(0\le c_i\le t\).  Since
\(t(t+1)^{d-1}<(t+1)^d\), some grid point gives a primitive
integer \(\alpha\).  Every conjugate of \(\alpha\) has absolute
value at most \(d^{d+3}\Delta\).  The elementary symmetric
formulae for its minimal polynomial give (B.A4).

For each polynomial, decide irreducibility, compute its integer ring
and discriminant, and retain it exactly when its field is Galois
and has no finite ramification outside \(\ell\).  All these are
finite exact operations.  For example, the integer ring can be
computed by enumerating the finitely many intermediate lattices
between \(\Z[\alpha]\) and its trace-dual lattice, checking
multiplication and integrality, and taking the maximal integral
order.  Polynomial factorization over number fields decides normality
and constructs composita and automorphisms.  These operations,
including their unconditional algorithmic form, are described by
Lenstra in \S\S2--4~\cite{Lenstra1992Algorithms}; no polynomial-time
claim is needed here.  Include \(\Q\) separately.

Let \(L\) be the compositum of all the retained fields and
\(\Q(\mu_q)\), and compute \(G=\operatorname{Gal}(L/\Q)\).
Choose a prime of \(L\) above \(\ell\) and let \(D_\ell\subset G\)
be its decomposition group.  It is computable by testing which
automorphisms preserve that prime ideal.  The action \(a_g\) is
read off from the action on the specified \(q\)-th root of unity.
For variables \((z_g)_{g\in G},b\) in \(\Z/q\), form the
following finite system of linear congruences:
\[
 \begin{aligned}
 z_{gh}-z_g-a_gz_h&=0 &&(g,h\in G),\\
 z_d-(a_d-1)b&=0 &&(d\in D_\ell).
 \end{aligned}
\tag{B.A5}
\]
Let \(Z_q\) be its solution module and quotient it by the cyclic
submodule generated by
\[
 \bigl((a_g-1)_{g\in G},1\bigr).
\tag{B.A6}
\]
The resulting finite abelian group is exactly
\(H^1_c(U,\Z/q(1-n))\).  Indeed, the mapping cone defining
compact support represents a degree-one cocycle as a global
crossed cocycle together with a local degree-zero trivialization;
its cocycle and boundary equations are (B.A5)--(B.A6).  At the
real place the modified Tate complex is acyclic because \(\ell\)
is odd.  Formula (B.A2) proves that every global cocycle occurs on
\(G\).  The local trivialization condition depends only on its
image \(D_\ell\), so no additional local cocycles have been
discarded.  Changing the chosen prime conjugates the presentation
and gives the canonical isomorphic compact-support group.

Finite-coefficient Artin--Verdier duality now gives
\[
 H^2(U,\Z/q(n))
 \simeq\operatorname{Hom}\bigl(Z_q/\langle\text{(B.A6)}\rangle,
                                  \Q/\Z\bigr).
\tag{B.A7}
\]
Smith normal form computes both cyclic factors and generators on
the right.  More explicitly, if \(x_j\) are the returned generators
of orders \(d_j\), the characters with value \(1/d_j\) on
\(x_j\) and zero on the other generators are the requested
dual cohomological generators.

Finally write \(C_{n,\ell}=H^2(U,\Z_\ell(n))\), which is finite
for \(n\ge2\).  The odd-primary Galois group here has
\(\ell\)-cohomological dimension two.  The finite coefficient
groups satisfy the Mittag--Leffler condition; hence continuous
cohomology in degree three is zero.  The coefficient exact sequence
therefore identifies the group in (B.A7) with
\(C_{n,\ell}/\ell^r C_{n,\ell}\).
Compute successively for \(r=1,2,\ldots\), and stop once two
successive orders agree.  If
\(C_{n,\ell}=\bigoplus_j\Z/\ell^{e_j}\), the orders have
exponents \(\sum_j\min(r,e_j)\).  Equality at \(r,r+1\)
holds exactly when every \(e_j\le r\).  Thus the stopping
criterion is sufficient and, by finiteness, is eventually reached.
At that stage (B.A7) is the whole group with its generators.

Positive-twist localization, or equivalently the prime-to-\(\ell\)
order \(\ell^{n-1}-1\) of the relevant finite-field term, identifies
\(C_n[\ell^\infty]\) with \(C_{n,\ell}\).  The finite prime
bound (B.A1) and the just-proved termination for each prime finish
the algorithm.  No enumeration in this proof is asserted to be
computationally practical, and the finite verification accompanying
it tests the cocycle presentation, not all of its astronomically
large number-field searches.
\end{proof}


\subsubsection{A terminating construction of primitive integral regulator classes}

\begin{theorem}[An integral generator specified by finite searches]
\label{thm:B-primitive-generator}
\status{Proved as an effective construction}
For each odd integer \(m\ge3\), the following procedure specifies an
higher Chow cycle \(\gamma_m\) over \(\Q\) with integer
coefficients whose canonical integral-base class generates
\[
 H^1_{\et}(\Spec\Z,\Z(m))/\mathrm{tors}
 \simeq\Z.
\]
Every search terminates.  The procedure requires the finite arithmetic
computation of \(C_m\) in Theorem~\ref{thm:B-effective-odd-torsion}
and Proposition~\ref{prop:B-two-primary}, but requires no uniform
vanishing conjecture for these groups.  Its output is a cycle given by
rational polynomial equations and integer multiplicities, together with
an explicitly found boundary certifying the saturation step.
\end{theorem}
\begin{proof}
We first specify exactly the finite tests used by our searches.  Put
\(\square=\mathbf P^1\setminus\{1\}\), and write
\(z^m(\Q,r)\) for the integral cubical higher Chow complex, modulo
degenerate cycles.  At stage \(H=1,2,\ldots\), enumerate at most \(H\) polynomial
generators per ideal, coordinate degrees at most \(H\), rational
coefficients with numerator of absolute value at most \(H\) and
denominator between one and \(H\), at most \(H\) cycle
components, and multiplicities with numerator and denominator bounded
by \(H\).  Order each finite stage lexicographically.  Every
finite rational cycle description eventually occurs.  Multihomogeneous
ideals in \((\mathbf P^1)^r\), with components contained in a
coordinate \(1\) removed, describe all the required cycles.  Ideal
factorization, dimension, saturation, and intersection multiplicities
compute their irreducible components and their restrictions to every
face \(z_i=0,\infty\).  These are finite operations over \(\Q\).
For ideal membership, dimension, primary decomposition and associated
primes, respectively, one can use Seidenberg's numbered statements
5, 6, 36 and 46; the remark following statement 46 verifies the
required effectiveness hypotheses over \(\Q\)~\cite{BSeidenberg1974}.
Intersection multiplicities are the lengths of the zero-dimensional
quotients at the computed generic points, computed by finite basis
linear algebra over their function fields.
The admissibility test is the finite set of dimension inequalities for
all faces.  Equality of two cycles is tested component by component;
degeneracy is tested by elimination in each coordinate and equality
with the inverse image of the resulting ideal.  Thus the differential
\(\partial\), the assertion \(\partial Z=0\), and any specified
boundary identity are decidable for each candidate.  The procedure
never asks whether an arbitrary cycle is a boundary without searching
for a particular bounding chain.

We also use rational coefficients during the first search.  A cycle is
called real-proper with exactly the finite family of conditions used
by the regulator complex.  For \(r=2m-1\), for each
\(J\subset\{1,\ldots,r\}\), each choice of a zero or infinity
face on \(J\), and each \(0\le k\le r\), put
\(I=\{1,\ldots,k\}\setminus J\).  Require proper real
intersection with the chosen faces and the cuts
\(z_i\in\R_{<0}\) for \(i\in I\).  These are precisely the
conditions specified in Kerr--Li, footnote~4~\cite{BKerrLi2020}.  After writing complex coordinates in real and
imaginary parts, these are finite dimension questions for semialgebraic
sets over \(\Q\).  Cylindrical algebraic decomposition decides them:
take a decomposition adapted to the finitely many polynomial equalities
and inequalities and read off the dimensions of its cells.  Collins's
Section 3, algorithms DECOMP and ELIM and Theorem 8, gives the
required finite procedure~\cite{BCollins1975}.  No
approximation to a period occurs in this test.

Here is the first search.  Set \(t_i=z_i/(z_i-1)\) and let \(Z_m\)
be the closure of the following graph in \(\square^{2m-1}\):
\[
 \left(
 \frac{t_1}{t_1-1},\ldots,\frac{t_{m-1}}{t_{m-1}-1},
 1+t_1\cdots t_{m-1},t_1^2,\ldots,t_{m-1}^2
 \right).
\]
For \(1\le i\le m-1\), let \(p_i\) project to the first
\(m-i\) coordinates, and put
\[
 f_a=1+\prod_{j=1}^a\frac{z_j}{z_j-1}.
\]
Enumerate tuples of rational admissible real-proper cycles
\((W_1,\ldots,W_{m-1})\) in \(z^m(\Q,2m-1)_\Q\)
until
\[
 \operatorname{supp}(W_i)\subset p_i^{-1}V(f_{m-i}),
 \qquad
 \partial\left(Z_m+\sum_iW_i\right)=0.
\tag{B.G1}
\]
Use the zero divisor in the open cube in this formula, so clearing
denominators does not add a component in a removed face.  These support
conditions are again ideal-membership tests.  Kerr--Yang's
Theorem~3.3 and Remark~3.4, specialized to \(N=2\), prove that a
tuple satisfying (B.G1) exists.  Their existence argument precedes and
does not use the explicit formulae in their \S4: Proposition~2.13
annihilates the residue class, and its bounding chain is decomposed
according to the faces of the relative cycle complex.  The pieces are
the \(W_i\), with exactly the support conditions (B.G1).
Consequently the first search terminates~\cite[Theorem~3.3 and Remark~3.4]{BKerrYang2018}.  Write
\(\widetilde Z_m=Z_m+\sum W_i\), let \(D_m\) be a common
denominator of its actual coefficients, and set
\[
 d_m=(m-1)!D_m,\qquad \xi_m=d_m\widetilde Z_m.
\tag{B.G2}
\]
This is an integral closed cycle.  Indeed, the cubical cycle groups
modulo degeneracies are free abelian on the nondegenerate irreducible
cycles, so a rationally zero integral boundary is already zero.

We check the regulator normalization before saturating.  In the
unscaled real coordinate of
\(H^1_D(\Spec\Q/\R,\R(m))\simeq\R\), Kerr--Yang's
Theorem~3.8 gives
\[
 \rho(\xi_m)
 =-d_m2^{m-1}\operatorname{Li}_m(-1)
 =d_m(2^{m-1}-1)\zeta(m)>0.
\tag{B.G3}
\]
One can verify the vanishing of the correction terms directly: on
their support the product of the corresponding \(t_j\)'s is
\(-1\), whereas on the simultaneous negative-real cuts all these
\(t_j\) lie in \([0,1]\).  The graph term is the iterated
integral in the proof of their Theorem~3.8.  Taking real parts removes
the ambiguity by \((2\pi i)^m\Q\), because \(m\) is odd.
We use this equation, not the inconsistent \(N=2\) sentence following
their Corollary~3.9~\cite[Theorem~3.8]{BKerrYang2018}.

Put \(c_m=|C_m|\), now an explicitly computed positive integer.
Let \(g\) be a positive primitive class in
\(H^1_{\et}(\Z,\Z(m))/\mathrm{tors}\).
Flach--Morin's regulator in \S5.8.3 uses the target lattice
\((2\pi i)^{m-1}\Z\), so its covolume is
\(R_m=\rho(g)/(2\pi)^{m-1}\).  Both maps are the Beilinson
regulator from the same motivic complex; passage from the complex
Deligne quotient to this real coordinate is the real-part projection,
with no factorial or additional factor two.  The factor
\(2/(2\pi i)^{m-1}\) in Kerr--Yang Lemma 2.8 converts their
Beilinson regulator to the Borel regulator; Flach--Morin uses the
Beilinson regulator itself, so that conversion is not made here.
Since \(w_m=2\),
their equation (88), proved for \(\Q\) by Proposition~5.35,
and the functional equation of \(\zeta\) give
\[
 |\zeta'(1-m)|
 =\frac{c_m\rho(g)}{2(2\pi)^{m-1}}
 =\frac{(m-1)!\zeta(m)}{2(2\pi)^{m-1}}.
\]
It follows that the index of \(\xi_m\) in the free lattice is
the explicitly known positive integer
\[
 h_m=\frac{d_m(2^{m-1}-1)c_m}{(m-1)!}
     =D_m(2^{m-1}-1)c_m.
\tag{B.G4}
\]
All factors and lattices here are those specified in the opened
source~\cite[\S5.8.3, equation~(88), Proposition~5.35]{BFlachMorin2018}.
The integer in (B.G4) is determined by exact arithmetic; no numerical
recognition of a regulator ratio is used.

For completeness the cycles may be searched over \(\Q\).
Positive-twist localization identifies
\[
 H^1_M(\Spec\Z,\Z(m))
 \xrightarrow{\ \sim\ }
 H^1_M(\Spec\Q,\Z(m))=CH^m(\Q,2m-1),
\tag{B.G5}
\]
because the adjacent finite-field groups are
\(H^{-1}_M(\F_p,\Z(m-1))\) and
\(H^0_M(\F_p,\Z(m-1))\), both zero for \(m-1>0\).
Geisser's Theorem~3.2 gives the higher-Chow localization sequence,
and his Theorem~1.2(2) gives the motivic-to-etale comparison through
degree \(m+1\), under the norm-residue theorem.  That theorem is
Voevodsky's Theorem~6.16, with the integral field comparison in
Theorem~6.18.  Thus the comparison invoked in Flach--Morin
\S5.8.3 is unconditional in the degree used here
\cite[Theorems~1.2(2), 3.2]{BGeisserDedekind2004}
\cite[Theorems~6.16, 6.18]{BVoevodsky2011}.
The torsion of the group in (B.G5) has order \(w_m=2\).

Finally enumerate pairs of integral admissible cycles
\[
 \gamma\in z^m(\Q,2m-1),\qquad
 \eta\in z^m(\Q,2m)
\]
until the finite tests give
\[
 \partial\gamma=0,
 \qquad \partial\eta=2(h_m\gamma-\xi_m).
\tag{B.G6}
\]
Equation (B.G4) proves existence: any integral representative of the
positive primitive free class satisfies
\(h_m[\gamma]-[\xi_m]\in\mathrm{tors}\), and twice this
difference is a boundary by the definition of higher Chow homology.
Hence this search also terminates.  Conversely (B.G6) implies that
\(h_m[\gamma]=[\xi_m]\) in the free quotient.  Since the latter
class has index exactly \(h_m\), \([\gamma]\) is primitive.
We take the first pair in the specified enumeration and define
\(\gamma_m=\gamma\).  The returned \(\eta\) certifies the
integral saturation.  This proves the claimed construction.
\end{proof}

\begin{lemma}[The degeneracy in the displayed uniform correction]
\label{lem:B-KY-degenerate}
\status{Proved source-formula defect; no general boundary formula asserted}
For every weight \(m\ge4\), the displayed \(W_2^{(i)}\)
terms in Kerr--Yang \S4D are cubically degenerate, or zero,
as rational higher Chow cycles.  Their existence theorem used above
does not depend on those displayed terms.
\end{lemma}
\begin{proof}
Every coordinate of a displayed \(W_2^{(i)}\), other than the
last coordinate \(v-1\), is homogeneous of degree zero in
\((u,v)\).  In the nontrivial ratio for \(i\le m-2\), the
numerator has degree \(m-3\) and the denominator degree
\(1+(m-4)=m-3\).  For the last term the denominator has
degree \(-1+(m-2)=m-3\).  The remaining coordinates are
ratios such as \(v t_j^N/u\), \(v/u\), or
\(u/(v\prod t_j^N)\), hence have the same homogeneity.
On the dense chart \(v\ne0\), substitute \(u=sv\).
All but the final coordinate now depend only on \(s,t_j\),
while \(v-1\) independently varies through a dense open subset
of the final cube.  The image closure is consequently the pullback
of its projection with that coordinate omitted, which is a
degenerate cycle.  If one of the other coordinates is identically
one, the cycle is already empty in the open cube.

The symbolic check for weights four, five and six, with \(N=2\),
is preserved in \texttt{source\_and\_graph\_checks.py}; the
homogeneity argument proves the assertion in every stated weight.
Kerr--Li independently records that the first weight-four formula
is not closed, because its proposed \(W_2\) correction is
degenerate~\cite[p.~548, footnote~10]{BKerrLi2020}.
We do not infer a complete uniform boundary computation from this
single defect.  The construction in
Theorem~\ref{thm:B-primitive-generator} instead uses the independent
existence theorem and tests the boundary of its actual output.
\end{proof}



\begin{proposition}[An exact Vandiver gate inside B.4]
\label{prop:B-Vandiver}
\status{Proved}
Fix an odd prime \(p\), set \(L=\Q(\zeta_p)\),
\(\Delta=\Gal(L/\Q)\), and let \(A\) be the \(p\)-primary
part of \(\BCl(\mathcal O_L)\).  Write \(A[j]\) for its
\(\omega^j\)-eigenspace, where \(\omega\) is the Teichmuller
character.  For every \(m\ge2\),
\[
 H^2_{\Bcts}(\Z[1/p],\Z_p(m))/p
 \simeq (A/p\otimes\mu_p^{\otimes(m-1)})^\Delta
 \simeq A[1-m]/p.\tag{B.23}
\]
Consequently the simultaneous assertion
\[
 C_m[p^\infty]=0\quad\text{for every odd }m\ge3
\]
is equivalent to the \(p\)-instance of Vandiver's conjecture,
\(p\nmid h(\Q(\zeta_p)^+)\).
\end{proposition}
\begin{proof}
There is exactly one prime of \(L\) over \(p\), and it is
principal.  Thus removing it does not change \(A\).  Since \(L\)
has no real places, the invariant-sum description of the Brauer group
gives \(\Br(\mathcal O_L[1/p])[p^\infty]=0\): a single
local invariant has zero sum only when it vanishes.  Kummer theory
therefore identifies
\(H^2(\mathcal O_L[1/p],\mu_p)\) with \(A/p\).
Tensoring by \(\mu_p^{\otimes(m-1)}\) and using the exactness
of \(\Delta\)-invariants, since \(|\Delta|=p-1\), gives the
middle term of (B.23).  The coefficient sequence for multiplication
by \(p\) has no obstruction in degree three: the \(p\)-cohomological
dimension of this number-ring Galois group is two.  This proves
(B.23).  The group on its left before reduction modulo \(p\) is
finite for \(m\ge2\); Nakayama's lemma therefore detects its
vanishing by (B.23).

Etale motivic comparison and localization identify its finite
\(p\)-primary group with \(C_m[p^\infty]\).  The local term at
\(p\) is zero here, because the finite-field positive-twist groups
in twist \(m-1\ge1\) have order prime to \(p\).  As odd
\(m\ge3\) runs through residue classes modulo \(p-1\),
\(1-m\) runs through every even residue class.  Their sum is the
plus part of \(A\); its vanishing is exactly the indicated
class-number condition, since the quadratic extension to the real
subfield has degree prime to \(p\).  This also rederives the precise
equivalence recorded in Kurihara's Corollary~1.5~\cite{BKurihara1992}.
\end{proof}

This reduction does not logically say that every possible algorithm
for individual \(C_m\) must first prove Vandiver.  It says precisely
that the often-proposed uniform answer ``zero in all these odd
twists'' is equivalent to Vandiver, and that an all-twist table must
settle or retain these eigenspaces.  The 2026 paper \emph{Almost all
primes are partially regular}, Theorem~1.1 and its partial Vandiver
corollary, treats a growing but restricted range of eigenspaces for
a density-one set of primes; its statement does not assert the
universal vanishing in Proposition~\ref{prop:B-Vandiver}
\cite[Theorem~1.1 and Corollary~1.4]{BChen2026}.

\begin{example}[A certified irregular-prime instance]
\label{ex:B-Vandiver37}
\status{Proved by a complete arithmetic certificate}
Let \(L=\Q(\zeta_{37}+\zeta_{37}^{-1})\), with
\(\zeta_{37}\) a primitive thirty-seventh root of unity.
The preserved PARI/GP computation constructs its degree-eighteen
real cyclotomic polynomial, returns the trivial computed class
group, and returns \(1\) from \texttt{bnfcertify(B,1)}.
The documented meaning of that flag is that the actual class
group is a quotient of the computed group; since the latter is
trivial, the actual class group is trivial
\cite{BPariClass2026}.  Consequently
\(C_m[37^\infty]=0\) for every odd \(m\ge3\), by
Proposition~\ref{prop:B-Vandiver}.

The same exact computation gives
\[
 \frac{B_{32}}{32}=-\frac{7709321041217}{16320},\qquad
 v_{37}(B_{32}/32)=1,
\]
so this is an irregular prime instance.  The executable input is
\path{part_b/code/vandiver37_quotient.gp}; the full polynomial,
certificate output, PARI version \(2.19.1\), and elapsed time
\(47\) milliseconds are in
\path{part_b/output/vandiver37_quotient.txt}.
The interrupted stronger full-unit certification is retained as
unused output and plays no role in this proof.
\end{example}
\paragraph{Localization and the zeta order.}
For a finite set \(S\) of rational primes, the rational absolute
localization triangle has support term
\[
 \bigoplus_{v\in S}\BRGamma_{\BW}(\F_v,\Q(n-1))[-2].
\]
By (B.1) this is zero if \(n\ne1\).  This proves the localization
isomorphism asserted in 15:173, with exactly the corrected exclusion
\(n\ne1\), for the specified rational absolute theory.  In twist
one, its two support groups give instead the unit and invariant-sum
sequence (B.17).  An identification with ordinary motivic cohomology
would destroy the degree-three fundamental class, so localization
must be taken within the same theory throughout.

The ranks in (B.5) do not by themselves give the order of
\(\zeta(s)\) by the naive weighted Euler characteristic: for
example the isolated line in twist three would give \(-1\), while
\(\zeta(3)\ne0\).  The needed archimedean and regulator contribution
is explicit in the compactly supported Weil--Arakelov theory.  For
\(X=\Spec\Z\), Flach--Morin's computation in \S5.8.3 gives
\[
 H^i_{{\rm ar},c}(X,\widetilde{\R}(n))=
 \begin{cases}
 \R,&n\le-2\text{ even},\ i=1,2,\\
 \R,&n=1,\ i=2,3,\\
 0,&\text{otherwise}.
 \end{cases}\tag{B.24}
\]
It follows, with order negative for a pole, that
\[
\sum_i(-1)^i i\dim_\R H^i_{{\rm ar},c}
=\operatorname{ord}_{s=n}\zeta(s)
=\begin{cases}1,&n\le-2\text{ even},\\-1,&n=1,\\0,&\text{otherwise}.
\end{cases}\tag{B.25}
\]
Here is why no unproved leading-value formula is being used: the
Borel regulator is an isomorphism after tensoring with \(\R\),
so the positive-twist regulator lines cancel against their real
targets in the defining cone.  For negative twists the dual lines
and the real-place term remain in the two degrees shown; at twist
one the degree map leaves the two fundamental real lines.  These
are the cone computations displayed on pp.~1505--1506 of the
source.  Independently, the functional equation, the simple pole at
one, and the nonvanishing Euler product for real \(s>1\) give
exactly the right-hand side of (B.25); \(\zeta(0)=-1/2\) covers
the remaining edge case.  Lead B2 is therefore correct about the
dual regulator lines and the negative even zeros, with this
archimedean completion required for the full order formula.

\subsection{B.5: finite duality and the rational statement that survives}

\begin{theorem}[Finite absolute duality, with the coefficient category specified]
\label{thm:B-finite-duality}
\status{Proved}
Let \(X\) be regular, separated, and of finite type over \(\Spec\Z\),
of pure absolute dimension \(d\).  Let \(F\) be a bounded
constructible complex of \(\Z/m\)-sheaves on \(X_{\et}\).
There is a perfect pairing of finite groups
\[
 \widehat H_c^i(X,F)\ \times\
\operatorname{Ext}^{2d+1-i}_{X,\Z/m}
       (F,\Z/m(d))\longrightarrow\Q/\Z.\tag{B.26}
\]
Here \(\Z/m(d)\) is the etale motivic complex, the Ext is in the
derived category of \(\Z/m\)-sheaves, and the hat denotes the
Tate modification at real places.  If \(m\) is invertible on \(X\),
\(\Z/m(d)\simeq\mu_m^{\otimes d}\).  Thus, for a locally
constant finite free \(F\), the second factor is the cohomology of
the usual twisted dual sheaf.  The real-place modification disappears
for odd \(m\), for schemes with no real fibre, and for schemes over
a finite field; properness over \(\Spec\Z\) alone does not remove it.
\end{theorem}
\begin{proof}
Geisser--Schmidt's Theorem~4.6, applied over the global field \(\Q\),
gives
\[
\widehat H_c^i(X,F)\times
\operatorname{Ext}^{2-i}_{X}(F,\Z_X^c(0))\longrightarrow\Q/\Z.
\]
Its hypotheses allow any bounded constructible torsion complex here;
the exclusion of characteristic-primary torsion in their positive
characteristic global-field formulation is irrelevant when the base
is \(\Spec\Z\).  For our regular pure-dimensional \(X\), the
cycle-complex indexing gives
\(\Z_X^c(0)/m=\Z/m(d)[2d]\).
Their Lemma~3.7 gives
\[
\BRHom_X(F,\Z_X^c(0))\simeq
\BRHom_{X,\Z/m}(F,\Z_X^c(0)/m)[-1].
\]
The two shifts change degree \(2-i\) to \(2d+1-i\), proving
(B.26).  Their Theorem~3.9 identifies the cycle complex with
\(\mu_m^{\otimes d}[2d]\) when \(m\) is invertible.  In the
smooth case this recovers precisely Milne's Theorem~II.7.6 and
Corollary~II.7.7, whose relative dimension \(d-1\) gives total
degree \(2(d-1)+3=2d+1\)~\cite{BGeisserSchmidt2018,AMilneADT}.
Milne's updated addendum explicitly directs the regular flat extension
to this result~\cite{BMilneErrata2025}.

The correction at infinity is part of the compact-support functor,
not of the algebraic properness assumption.  It is the cone of
restriction to Tate cohomology for \(\Gal(\C/\R)\), as in
Geisser--Schmidt's equations~(2.12)--(2.13).  Tate cohomology of that
group vanishes after inverting two, but for the trivial module
\(\Z/2\) it is \(\Z/2\) in every degree.  This proves both
the stated exceptions and the need to retain the correction in general.
\end{proof}

There is a concrete characteristic-primary check on the notation.
For a smooth complete curve \(C/\F_q\), \(q\) a power of \(p\),
one has
\(\Z/p^a(1)=\nu_a^1[-1]\), where \(\nu_a^1\) is the sheaf
of logarithmic de Rham--Witt one-forms.  Hence the dualizing complex
is \(\nu_a^1[2]\).  Milne's Theorem~II.7.11 gives the perfect
pairing
\[
 H^i(C,\Z/p^a)\times H^{2-i}(C,\nu_a^1)
 \longrightarrow\Z/p^a.
\]
It is exactly (B.26) after the coefficient shift.  Replacing
\(\Z/p^a(1)\) by the etale sheaf \(\mu_{p^a}\), which is zero
on a smooth reduced characteristic-\(p\) scheme, would instead give
a false assertion.  This also checks the edge case \(p=2\).

\begin{proposition}[Finite duality without separatedness]
\label{prop:B-finite-duality-nonsep}
\status{Proved for compact support defined by open-cover descent}
The finite pairing (B.26) extends to every regular finite-type
\(\Z\)-scheme of pure dimension \(d\), when compact support
on a possibly nonseparated scheme is defined by the following
finite open-cover construction.  This definition agrees with the
usual Tate-modified compact support when the scheme is separated.
\end{proposition}
\begin{proof}
Choose a finite affine open cover \(X=\bigcup_{j=1}^N U_j\).
Every nonempty intersection \(U_I=\bigcap_{j\in I}U_j\) is
quasicompact and separated: it is an open subscheme of any one
of its affine members.  It is therefore of finite type over
\(\Spec\Z\), and Theorem~\ref{thm:B-finite-duality} applies
to it.  Use functorial cochain models, including the real-place
cones, as in Geisser--Schmidt \S2; totalization is therefore of
an actual diagram of complexes~\cite{BGeisserSchmidt2018}.  Define
\[
 \begin{gathered}
 \widehat R\Gamma_c(X,F):=\operatorname{Tot} C_c^\bullet(\mathcal U,F),\\
 C_c^{-r}(\mathcal U,F):=
 \bigoplus_{|I|=r+1}\widehat R\Gamma_c(U_I,F),\qquad 0\le r<N.
 \end{gathered}
\]
where the notation means the finite alternating Cech total complex,
with the \(|I|=r+1\) term in Cech degree \(-r\), and the
structure maps are extension by zero for open immersions.
The real-place Tate modification is included in
each term; its maps are natural for these open immersions.

This construction is independent of the cover.  For two covers
use their intersection refinement.  On each separated intersection,
the usual compact-support Mayer--Vietoris sequence identifies the
refinement totalization with the original term.  Totalizing the
finite double Cech complex in either order proves both refinement
isomorphisms and their compatibility.  The same argument for a
separated \(X\) identifies the definition with its usual compact
support.  No compactification of a nonseparated scheme is assumed.

For the duality, put \(D_X=\Z/m(d)\).  The finite alternating
resolution of \(F\) by the sheaves \(j_{I!}(F|_{U_I})\)
is exact, as is checked on each geometric stalk by the augmented
simplex of the members containing that point.  Open restriction
of the motivic cycle complex gives \(D_X|_{U_I}=D_{U_I}\),
and adjunction gives
\[
 R\operatorname{Hom}_{X,\Z/m}
 (j_{I!}(F|_{U_I}),D_X)
 \simeq R\operatorname{Hom}_{U_I,\Z/m}(F|_{U_I},D_{U_I}).
\]
Thus the dual side is the corresponding finite Cech cochain
totalization.  The local trace pairings of (B.26) are compatible
with the open-immersion maps, so their duality isomorphisms
identify these two totalizations with the shift \(2d+1\).
All totalizations here use finitely many terms.  Their groups in
each degree are finite, and \(\operatorname{Hom}(-,\Q/\Z)\)
is exact.  Taking cohomology therefore gives the perfect pairing
(B.26) on \(X\), with precisely the declared compact-support
convention.
\end{proof}

\begin{theorem}[Rational duality for the Tate cases proved here]
\label{thm:B-rational-duality}
\status{Proved}
For \(X=\Spec\mathcal O_K\), the Flach--Morin compactified
complexes have perfect rational pairings
\[
 H^i_{\BW}(\bar X,\Q(n))\times
 H^{3-i}_{\BW}(\bar X,\Q(1-n))\longrightarrow\Q.
 \tag{B.27}
\]
For \(U=\Spec\mathcal O_K[1/S]\), the same assertion pairs
the localized rational absolute complex with the compact-support
complex, with the archimedean completion inherited from \(\bar X\).
For a smooth projective geometrically connected curve \(C/\F_q\),
the corresponding Weil--etale pairing is
\[
 H^i_{\BW}(C,\Q(n))\times H^{3-i}_{\BW}(C,\Q(1-n))
 \longrightarrow\Q.
 \tag{B.28}
\]
Both statements hold for every integral twist \(n\), including
negative twists, and for finite direct sums of these Tate objects.
\end{theorem}
\begin{proof}
For the proper number ring, all the finite-generation and duality
conditions used in Flach--Morin's construction hold in dimension
one, as verified in their \S5.8.3.  Their Theorem~3.22 therefore
applies, and Corollary~3.12 identifies its rational pairing with
cross-evaluation between ordinary motivic cohomology and the shifted
dual of the complementary twist.  This proves (B.27) and explains
why ordinary motivic cohomology alone lacks its degree-three partner.

For the open scheme put
\(D_S(n)=\bigoplus_{v\in S}\BRGamma_{\BW}(\F_v,\Q(n))\).
Compact support is the fibre of restriction
\(\BRGamma_{\BW}(\bar X,\Q(n))\to D_S(n)\).
Localization is the cone of the Gysin map
\(D_S(n-1)[-2]\to\BRGamma_{\BW}(\bar X,\Q(n))\).
The point calculation (B.1) gives a perfect trace pairing between
\(D_S(n)\) and \(D_S(-n)\) of total degree one.  Local
restriction and the Gysin map are adjoint under that pairing and
the global trace: this follows already for finite coefficients from
(B.26), and on the rational point summands is the equality between
degree and the invariant-sum map.  Dualizing the compact-support
triangle in complementary twist therefore gives exactly the
localization triangle.  Since evaluation is perfect on its two
known terms, the induced cone pairing is perfect as well.  This is
the asserted open-scheme duality, and specifies which compact
support is meant.

For the projective curve, Lichtenbaum's Theorem~5.4 gives
\[
 \BRGamma_{\BW}(C,\Gm)\simeq
 \BRHom(\BRGamma_{\BW}(C,\Z),\Z[-2]).
\]
Since \(\Z(1)=\Gm[-1]\), this is the integral source of
the degree-three pairing in (B.28) for twists zero and one.
Their rational groups are \(\Q\) in degrees \(0,1\) at twist
zero and \(\Q\) in degrees \(2,3\) at twist one, by
Lichtenbaum's Theorem~6.1.  For every other twist, the unconditional
curve comparison in Geisser's Theorem~9.5 and the Frobenius weights
make both sides zero.  This proves every case of (B.28).
\end{proof}

\begin{corollary}[Finite-monodromy coefficient motives]
\label{cor:B-Artin-Tate-duality}
\status{Proved}
Let \(X\) be either an open number ring or a smooth projective
curve over a finite field.  Let \(W\) be a finite-dimensional
rational representation of a finite quotient of
\(\pi_1^{\et}(X)\).  Then the rational absolute or Weil--etale
theory just specified has a common rational structure for all
\(W(n)\), every integer \(n\), and perfect pairings
\[
 H^i_{\rm abs}(X,W(n))\ \times\
 H^{3-i}_{\rm abs,c}(X,W^\vee(1-n))\longrightarrow\Q.
\]
Compact support includes the stated archimedean completion on a
number ring; it is unnecessary for a proper finite-field curve.
The degree-one realization is the Bloch--Kato finite-condition
group with the same geometric conditions at the omitted places.
\end{corollary}
\begin{proof}
Choose a finite etale Galois cover \(Y\to X\), with group \(G\),
trivializing \(W\).  Apply the preceding Tate calculations to
each connected component of \(Y\), over its number field or full
constant finite field.  Descent defines the rational complex by
\[
 R\Gamma_{\rm abs}(X,W(n))=
  \bigl(R\Gamma_{\rm abs}(Y,\Q(n))\otimes_\Q W\bigr)^G,
\]
and similarly with compact support.  The invariants functor is
exact, with projector \(|G|^{-1}\sum_{g\in G}g\), so it commutes
with scalar extension to every \(\Q_\ell\).  Restriction and
corestriction preserve the unramified, crystalline and de Rham
local kernels defining \(f\) and \(g\); their composite is
the extension degree.  At a retained place above \(\ell\),
the finite etale cover is unramified, so crystalline descent
applies.  At a removed place above \(\ell\), the condition is
\(g\), for which de Rham descent allows a ramified extension.
Thus the degree-one descent is exactly
the claimed Bloch--Kato group.  Higher comparison splittings
can be averaged over \(G\), giving the same graded assertion.

Combine the perfect pairing on \(Y\) with evaluation
\(W\otimes W^\vee\to\Q\) and divide the trace by \(|G|\).
Its restriction to invariant subspaces is perfect: for an invariant
nonzero vector, average any vector pairing nontrivially with it.
This proves the assertion and its compatibility with localization.
\end{proof}

\paragraph{The scope still to be proved.}
For a general pure motive \(M\), a rational analogue of (B.26)
must specify an absolute or Weil--etale motivic category, finite
local conditions, a trace, and a compact-support operation including
infinity.  The proved finite theorem does not create those rational
objects: tensoring a finite torsion pairing with \(\Q\) gives zero.
For a dualizable coefficient motive on a regular base of absolute
dimension \(d\), the corrected derived assertion is
\[
 R\Gamma_{\rm abs}(X,M)\simeq
 R\operatorname{Hom}_{\Q}
 \bigl(R\Gamma_{\rm abs,c}(X,M^\vee(d)),\Q\bigr)[-2d-1],
\]
with the compactified arithmetic convention just proved in dimension
one.  A compatible rational trace and realization into a perfect
local/global duality imply that this pairing is perfect, since
extension of scalars from \(\Q\) to any \(\Q_\ell\) is
faithfully flat.  The finite theorem alone supplies neither the
rational trace nor that realization.
Even on a number ring the equality of the expected degree-one
cohomology of \(V_\ell A\) with \(A(K)\otimes\Q_\ell\)
contains the Sha-finiteness assertion in (B.21).  Thus the arbitrary
pure-motive portion of B.5 is not proved by the all-twist Tate
calculation, and remains explicitly recorded as a separate task.

\begin{proposition}[The exact number-ring rational gate for a dual pair]
\label{prop:B-general-rational-gate}
\status{Proved equivalence for dual completions; the realization hypothesis is conjectural}
Fix a pure coefficient motive \(M/K\), and write
\(M^*=M^\vee(1)\).  Suppose rational motivic spaces
\(A^a(M)\) and \(A^a(M^*)\), for \(a=0,1\), are specified
with their realization maps.  They admit a finite-dimensional
rational dual completion having the graded Selmer comparisons
in every degree if and only if the degree-zero and degree-one
maps for both motives satisfy Fontaine's f-admissibility
conditions \(c^0_\ell,c^1_\ell\) for every \(\ell\).
For the prescribed motivic spaces, this is a genuine realization
conjecture, not the assertion that some abstract vector spaces of
the same dimensions can be chosen.
\end{proposition}
\begin{proof}
Under those conditions put
\[
\begin{aligned}
 A^2(M)&=A^1(M^*)^\vee,&
 A^3(M)&=A^0(M^*)^\vee,\\
 R A(M)&=\bigoplus_{i=0}^3 A^i(M)[-i].
\end{aligned}
\]
Evaluation, using \((M^*)^*=M\), gives the perfect pairing
\(RA(M)\otimes RA(M^*)\to\Q[-3]\).
The degree-zero and degree-one realizations are the assumed
maps.  The Poitou--Tate calculation (B.16), with the splittings
already specified, gives the degree-two and degree-three
comparisons.  Hence this is a rational dual completion with all
the required graded comparisons on the proper number ring.
Conversely any such completion, with the prescribed degree-zero
and degree-one maps, implies exactly those f-admissibility
conditions upon taking those two cohomology groups.  No existence
of a six-operation motivic category or canonical product on all
these completions has been deduced.  The proposition identifies
the exact remaining comparison gate for this specified graded
construction, rather than concealing it in higher-degree notation.
\end{proof}
